% Môn: Vật lý
% Lớp: 12
% Chương: Chương II. Khí lí tưởng
% Bài: L12C2 Bài 8. Mô hình động học phân tử chất khí
% Số câu: 162
% App đọc file này trực tiếp — sửa rồi Commit trên GitHub, bấm Đồng bộ GitHub .tex

% L12C2 Bài 8. Mô hình động học phân tử chất khí

% ===== Câu 1 =====
% ID: L12C2B8-01-DS
% Mức: TH
\dangbt{Chuyển động Brown và bằng chứng về chuyển động phân tử}
\begin{ex}
% ID: L12C2B8-01-DS
% Mức: TH
Xét các phát biểu sau trong dạng “Chuyển động Brown và bằng chứng về chuyển động phân tử”.
\choiceTF
{\True Theo mô hình động học phân tử, các phân tử khí — phát biểu đúng là: chuyển động hỗn loạn không ngừng.}
{Theo mô hình động học phân tử, các phân tử khí — phát biểu đứng yên tại vị trí cân bằng là đúng.}
{\True Kích thước phân tử so với khoảng cách trung bình giữa chúng trong khí lí tưởng được coi là — phát biểu đúng là: rất nhỏ.}
{Kích thước phân tử so với khoảng cách trung bình giữa chúng trong khí lí tưởng được coi là — phát biểu bằng nhau là đúng.}
\loigiai{
\begin{itemchoice}
\item (Đúng) Theo mô hình động học phân tử, các phân tử khí — phát biểu đúng là: chuyển động hỗn loạn không ngừng. (Vì): Phân tử khí luôn chuyển động nhiệt hỗn loạn không ngừng. b) Sai: nội dung không phù hợp. c) Đúng: Trong mô hình khí lí tưởng, kích thước phân tử được bỏ qua so với khoảng cách giữa chúng. d) Sai: nội dung không phù hợp
\item (Sai) Theo mô hình động học phân tử, các phân tử khí — phát biểu đứng yên tại vị trí cân bằng là đúng.
\item (Đúng) Kích thước phân tử so với khoảng cách trung bình giữa chúng trong khí lí tưởng được coi là — phát biểu đúng là: rất nhỏ.
\item (Sai) Kích thước phân tử so với khoảng cách trung bình giữa chúng trong khí lí tưởng được coi là — phát biểu bằng nhau là đúng.
\end{itemchoice}
}
\end{ex}

% ===== Câu 2 =====
% ID: L12C2B8-01-TL
% Mức: TH
\begin{ex}
% ID: L12C2B8-01-TL
% Mức: TH
Trong dạng “Chuyển động Brown và bằng chứng về chuyển động phân tử”, Trình bày các luận điểm chính của mô hình động học phân tử chất khí.
\choiceTF

\loigiai{
Chất khí gồm vô số phân tử rất nhỏ, chuyển động hỗn loạn không ngừng; khoảng cách giữa chúng lớn, tương tác ngoài va chạm không đáng kể; va chạm được coi là đàn hồi.
}
\end{ex}

% ===== Câu 3 =====
% ID: L12C2B8-01-TLN
% Mức: TH
\begin{ex}
% ID: L12C2B8-01-TLN
% Mức: TH
Trong dạng “Chuyển động Brown và bằng chứng về chuyển động phân tử”, Một bình chứa $6\times10^{23}$ phân tử trong thể tích $0,02\,m^3$. Tính mật độ phân tử theo đơn vị $10^{25}\,m^{-3}$.
\shortans{3}
\loigiai{
$n_0=N/V=6\times10^{23}/0,02=3\times10^{25}\,m^{-3}$.
}
\end{ex}

% ===== Câu 4 =====
% ID: L12C2B8-01-TN
% Mức: NB
\begin{ex}
% ID: L12C2B8-01-TN
% Mức: NB
Trong dạng “Chuyển động Brown và bằng chứng về chuyển động phân tử”, Theo mô hình động học phân tử, các phân tử khí
\choice
{\True chuyển động hỗn loạn không ngừng}
{đứng yên tại vị trí cân bằng}
{chỉ chuyển động khi được nung nóng}
{chuyển động cùng một hướng}
\loigiai{
Đáp án A. Phân tử khí luôn chuyển động nhiệt hỗn loạn không ngừng.
}
\end{ex}

% ===== Câu 5 =====
% ID: L12C2B8-02-DS
% Mức: TH
\begin{ex}
% ID: L12C2B8-02-DS
% Mức: TH
Xét các phát biểu sau trong dạng “Chuyển động Brown và bằng chứng về chuyển động phân tử”.
\choiceTF
{\True Lực tương tác giữa các phân tử khí lí tưởng được coi là — phát biểu đúng là: không đáng kể trừ khi va chạm.}
{Lực tương tác giữa các phân tử khí lí tưởng được coi là — phát biểu luôn rất lớn là đúng.}
{\True Chuyển động Brown là chuyển động — phát biểu đúng là: hỗn loạn của hạt nhỏ lơ lửng do va chạm phân tử.}
{Chuyển động Brown là chuyển động — phát biểu rơi tự do của hạt bụi là đúng.}
\loigiai{
\begin{itemchoice}
\item (Đúng) Lực tương tác giữa các phân tử khí lí tưởng được coi là — phát biểu đúng là: không đáng kể trừ khi va chạm. (Vì): Ngoài thời gian va chạm rất ngắn, tương tác phân tử được bỏ qua. b) Sai: nội dung không phù hợp. c) Đúng: Các phân tử môi trường va chạm không cân bằng làm hạt nhỏ chuyển động Brown. d) Sai: nội dung không phù hợp
\item (Sai) Lực tương tác giữa các phân tử khí lí tưởng được coi là — phát biểu luôn rất lớn là đúng.
\item (Đúng) Chuyển động Brown là chuyển động — phát biểu đúng là: hỗn loạn của hạt nhỏ lơ lửng do va chạm phân tử.
\item (Sai) Chuyển động Brown là chuyển động — phát biểu rơi tự do của hạt bụi là đúng.
\end{itemchoice}
}
\end{ex}

% ===== Câu 6 =====
% ID: L12C2B8-02-TL
% Mức: TH
\begin{ex}
% ID: L12C2B8-02-TL
% Mức: TH
Trong dạng “Chuyển động Brown và bằng chứng về chuyển động phân tử”, Giải thích chuyển động Brown.
\choiceTF

\loigiai{
Phân tử môi trường chuyển động hỗn loạn va chạm không cân bằng từ mọi phía vào hạt nhỏ, làm hạt chuyển động ngẫu nhiên.
}
\end{ex}

% ===== Câu 7 =====
% ID: L12C2B8-02-TLN
% Mức: TH
\begin{ex}
% ID: L12C2B8-02-TLN
% Mức: TH
Trong dạng “Chuyển động Brown và bằng chứng về chuyển động phân tử”, Một lượng khí có thể tích 4 lít được nén còn 2 lít. Mật độ phân tử tăng bao nhiêu lần?
\shortans{2}
\loigiai{
Số phân tử không đổi nên mật độ tỉ lệ nghịch thể tích: $n_2/n_1=V_1/V_2=2$.
}
\end{ex}

% ===== Câu 8 =====
% ID: L12C2B8-02-TN
% Mức: NB
\begin{ex}
% ID: L12C2B8-02-TN
% Mức: NB
Trong dạng “Chuyển động Brown và bằng chứng về chuyển động phân tử”, Kích thước phân tử so với khoảng cách trung bình giữa chúng trong khí lí tưởng được coi là
\choice
{\True rất nhỏ}
{rất lớn}
{bằng nhau}
{không xác định}
\loigiai{
Đáp án A. Trong mô hình khí lí tưởng, kích thước phân tử được bỏ qua so với khoảng cách giữa chúng.
}
\end{ex}

% ===== Câu 9 =====
% ID: L12C2B8-03-DS
% Mức: VD
\begin{ex}
% ID: L12C2B8-03-DS
% Mức: VD
Xét các phát biểu sau trong dạng “Chuyển động Brown và bằng chứng về chuyển động phân tử”.
\choiceTF
{\True Áp suất khí lên thành bình sinh ra chủ yếu do — phát biểu đúng là: phân tử va chạm thành bình.}
{Áp suất khí lên thành bình sinh ra chủ yếu do — phát biểu trọng lượng bình là đúng.}
{\True Khi nhiệt độ khí tăng, chuyển động nhiệt của phân tử nhìn chung — phát biểu đúng là: nhanh hơn.}
{Khi nhiệt độ khí tăng, chuyển động nhiệt của phân tử nhìn chung — phát biểu không đổi là đúng.}
\loigiai{
\begin{itemchoice}
\item (Đúng) Áp suất khí lên thành bình sinh ra chủ yếu do — phát biểu đúng là: phân tử va chạm thành bình. (Vì): Xung lượng truyền trong các va chạm phân tử với thành bình tạo áp suất. b) Sai: nội dung không phù hợp. c) Đúng: Nhiệt độ tăng làm động năng tịnh tiến trung bình tăng. d) Sai: nội dung không phù hợp
\item (Sai) Áp suất khí lên thành bình sinh ra chủ yếu do — phát biểu trọng lượng bình là đúng.
\item (Đúng) Khi nhiệt độ khí tăng, chuyển động nhiệt của phân tử nhìn chung — phát biểu đúng là: nhanh hơn.
\item (Sai) Khi nhiệt độ khí tăng, chuyển động nhiệt của phân tử nhìn chung — phát biểu không đổi là đúng.
\end{itemchoice}
}
\end{ex}

% ===== Câu 10 =====
% ID: L12C2B8-03-TL
% Mức: VD
\begin{ex}
% ID: L12C2B8-03-TL
% Mức: VD
Trong dạng “Chuyển động Brown và bằng chứng về chuyển động phân tử”, Giải thích vì sao chất khí gây áp suất lên thành bình.
\choiceTF

\loigiai{
Mỗi va chạm làm phân tử đổi động lượng và truyền xung lượng cho thành; tổng hợp rất nhiều va chạm tạo lực trung bình trên một đơn vị diện tích, tức áp suất.
}
\end{ex}

% ===== Câu 11 =====
% ID: L12C2B8-03-TLN
% Mức: VD
\begin{ex}
% ID: L12C2B8-03-TLN
% Mức: VD
Trong dạng “Chuyển động Brown và bằng chứng về chuyển động phân tử”, Nhiệt độ tăng từ 300 $K$ lên 600 K. Động năng tịnh tiến trung bình tăng bao nhiêu lần?
\shortans{2}
\loigiai{
Động năng trung bình tỉ lệ với nhiệt độ tuyệt đối nên tăng $600/300=2$ lần.
}
\end{ex}

% ===== Câu 12 =====
% ID: L12C2B8-03-TN
% Mức: TH
\begin{ex}
% ID: L12C2B8-03-TN
% Mức: TH
Trong dạng “Chuyển động Brown và bằng chứng về chuyển động phân tử”, Lực tương tác giữa các phân tử khí lí tưởng được coi là
\choice
{\True không đáng kể trừ khi va chạm}
{luôn rất lớn}
{lực hút không đổi}
{bằng trọng lực}
\loigiai{
Đáp án A. Ngoài thời gian va chạm rất ngắn, tương tác phân tử được bỏ qua.
}
\end{ex}

% ===== Câu 13 =====
% ID: L12C2B8-04-DS
% Mức: VD
\begin{ex}
% ID: L12C2B8-04-DS
% Mức: VD
Xét các phát biểu sau trong dạng “Chuyển động Brown và bằng chứng về chuyển động phân tử”.
\choiceTF
{\True Ba thông số trạng thái cơ bản của một lượng khí là — phát biểu đúng là: áp suất, thể tích, nhiệt độ.}
{Ba thông số trạng thái cơ bản của một lượng khí là — phát biểu khối lượng, trọng lượng, nhiệt lượng là đúng.}
{\True Va chạm giữa các phân tử khí lí tưởng được xem là — phát biểu đúng là: đàn hồi.}
{Va chạm giữa các phân tử khí lí tưởng được xem là — phát biểu không bảo toàn động lượng là đúng.}
\loigiai{
\begin{itemchoice}
\item (Đúng) Ba thông số trạng thái cơ bản của một lượng khí là — phát biểu đúng là: áp suất, thể tích, nhiệt độ. (Vì): Trạng thái lượng khí xác định bởi p, $V$, T. b) Sai: nội dung không phù hợp. c) Đúng: Mô hình khí lí tưởng coi các va chạm là đàn hồi. d) Sai: nội dung không phù hợp
\item (Sai) Ba thông số trạng thái cơ bản của một lượng khí là — phát biểu khối lượng, trọng lượng, nhiệt lượng là đúng.
\item (Đúng) Va chạm giữa các phân tử khí lí tưởng được xem là — phát biểu đúng là: đàn hồi.
\item (Sai) Va chạm giữa các phân tử khí lí tưởng được xem là — phát biểu không bảo toàn động lượng là đúng.
\end{itemchoice}
}
\end{ex}

% ===== Câu 14 =====
% ID: L12C2B8-04-TL
% Mức: VD
\begin{ex}
% ID: L12C2B8-04-TL
% Mức: VD
Trong dạng “Chuyển động Brown và bằng chứng về chuyển động phân tử”, So sánh chuyển động phân tử khi nhiệt độ tăng.
\choiceTF

\loigiai{
Nhiệt độ tuyệt đối tăng làm động năng tịnh tiến trung bình và tốc độ đặc trưng của phân tử tăng.
}
\end{ex}

% ===== Câu 15 =====
% ID: L12C2B8-04-TLN
% Mức: VD
\begin{ex}
% ID: L12C2B8-04-TLN
% Mức: VD
Trong dạng “Chuyển động Brown và bằng chứng về chuyển động phân tử”, Khí có $N=1,2\times10^{24}$ phân tử và mật độ $3\times10^{25}\,m^{-3}$. Tính thể tích theo lít.
\shortans{40}
\loigiai{
$V=N/n_0=0,04\,m^3=40$ lít.
}
\end{ex}

% ===== Câu 16 =====
% ID: L12C2B8-04-TN
% Mức: TH
\begin{ex}
% ID: L12C2B8-04-TN
% Mức: TH
Trong dạng “Chuyển động Brown và bằng chứng về chuyển động phân tử”, Chuyển động Brown là chuyển động
\choice
{\True hỗn loạn của hạt nhỏ lơ lửng do va chạm phân tử}
{thẳng đều của phân tử khí}
{rơi tự do của hạt bụi}
{tuần hoàn của chất lỏng}
\loigiai{
Đáp án A. Các phân tử môi trường va chạm không cân bằng làm hạt nhỏ chuyển động Brown.
}
\end{ex}

% ===== Câu 17 =====
% ID: L12C2B8-05-DS
% Mức: VDC
\begin{ex}
% ID: L12C2B8-05-DS
% Mức: VDC
Xét các phát biểu sau trong dạng “Chuyển động Brown và bằng chứng về chuyển động phân tử”.
\choiceTF
{\True Khi nén khí ở nhiệt độ không đổi, mật độ phân tử — phát biểu đúng là: tăng.}
{Khi nén khí ở nhiệt độ không đổi, mật độ phân tử — phát biểu giảm là đúng.}
{\True Hiện tượng khuếch tán chứng tỏ các phân tử — phát biểu đúng là: chuyển động không ngừng.}
{Hiện tượng khuếch tán chứng tỏ các phân tử — phát biểu đứng yên là đúng.}
\loigiai{
\begin{itemchoice}
\item (Đúng) Khi nén khí ở nhiệt độ không đổi, mật độ phân tử — phát biểu đúng là: tăng. (Vì): Cùng số phân tử trong thể tích nhỏ hơn nên mật độ phân tử tăng. b) Sai: nội dung không phù hợp. c) Đúng: Khuếch tán là hệ quả trực tiếp của chuyển động nhiệt hỗn loạn. d) Sai: nội dung không phù hợp
\item (Sai) Khi nén khí ở nhiệt độ không đổi, mật độ phân tử — phát biểu giảm là đúng.
\item (Đúng) Hiện tượng khuếch tán chứng tỏ các phân tử — phát biểu đúng là: chuyển động không ngừng.
\item (Sai) Hiện tượng khuếch tán chứng tỏ các phân tử — phát biểu đứng yên là đúng.
\end{itemchoice}
}
\end{ex}

% ===== Câu 18 =====
% ID: L12C2B8-05-TL
% Mức: VD
\begin{ex}
% ID: L12C2B8-05-TL
% Mức: VD
Trong dạng “Chuyển động Brown và bằng chứng về chuyển động phân tử”, Giải thích hiện tượng mùi nước hoa lan khắp phòng.
\choiceTF

\loigiai{
Phân tử nước hoa bay hơi rồi chuyển động nhiệt hỗn loạn, va chạm và khuếch tán từ nơi nồng độ cao sang thấp.
}
\end{ex}

% ===== Câu 19 =====
% ID: L12C2B8-05-TLN
% Mức: VD
\begin{ex}
% ID: L12C2B8-05-TLN
% Mức: VD
Trong dạng “Chuyển động Brown và bằng chứng về chuyển động phân tử”, Có 2 mol khí. Lấy $N_A=6,02\times10^{23}$. Số phân tử theo đơn vị $10^{23}$ là bao nhiêu?
\shortans{12,04}
\loigiai{
$N=2N_A=12,04\times10^{23}$.
}
\end{ex}

% ===== Câu 20 =====
% ID: L12C2B8-05-TN
% Mức: TH
\begin{ex}
% ID: L12C2B8-05-TN
% Mức: TH
Trong dạng “Chuyển động Brown và bằng chứng về chuyển động phân tử”, Áp suất khí lên thành bình sinh ra chủ yếu do
\choice
{\True phân tử va chạm thành bình}
{trọng lượng bình}
{lực hút của Trái Đất}
{thành bình nóng lên}
\loigiai{
Đáp án A. Xung lượng truyền trong các va chạm phân tử với thành bình tạo áp suất.
}
\end{ex}

% ===== Câu 21 =====
% ID: L12C2B8-06-DS
% Mức: TH
\dangbt{Các luận điểm cơ bản của mô hình động học phân tử chất khí}
\begin{ex}
% ID: L12C2B8-06-DS
% Mức: TH
Xét các phát biểu sau trong dạng “Các luận điểm cơ bản của mô hình động học phân tử chất khí”.
\choiceTF
{\True Theo mô hình động học phân tử, các phân tử khí — phát biểu đúng là: chuyển động hỗn loạn không ngừng.}
{Theo mô hình động học phân tử, các phân tử khí — phát biểu đứng yên tại vị trí cân bằng là đúng.}
{\True Kích thước phân tử so với khoảng cách trung bình giữa chúng trong khí lí tưởng được coi là — phát biểu đúng là: rất nhỏ.}
{Kích thước phân tử so với khoảng cách trung bình giữa chúng trong khí lí tưởng được coi là — phát biểu bằng nhau là đúng.}
\loigiai{
\begin{itemchoice}
\item (Đúng) Theo mô hình động học phân tử, các phân tử khí — phát biểu đúng là: chuyển động hỗn loạn không ngừng. (Vì): Phân tử khí luôn chuyển động nhiệt hỗn loạn không ngừng. b) Sai: nội dung không phù hợp. c) Đúng: Trong mô hình khí lí tưởng, kích thước phân tử được bỏ qua so với khoảng cách giữa chúng. d) Sai: nội dung không phù hợp
\item (Sai) Theo mô hình động học phân tử, các phân tử khí — phát biểu đứng yên tại vị trí cân bằng là đúng.
\item (Đúng) Kích thước phân tử so với khoảng cách trung bình giữa chúng trong khí lí tưởng được coi là — phát biểu đúng là: rất nhỏ.
\item (Sai) Kích thước phân tử so với khoảng cách trung bình giữa chúng trong khí lí tưởng được coi là — phát biểu bằng nhau là đúng.
\end{itemchoice}
}
\end{ex}

% ===== Câu 22 =====
% ID: L12C2B8-06-TL
% Mức: VDC
\dangbt{Chuyển động Brown và bằng chứng về chuyển động phân tử}
\begin{ex}
% ID: L12C2B8-06-TL
% Mức: VDC
Trong dạng “Chuyển động Brown và bằng chứng về chuyển động phân tử”, Nêu giới hạn của mô hình khí lí tưởng.
\choiceTF

\loigiai{
Mô hình phù hợp khi khí loãng, áp suất không quá cao và nhiệt độ không gần vùng hóa lỏng; khi đó thể tích riêng và tương tác phân tử có thể bỏ qua.
}
\end{ex}

% ===== Câu 23 =====
% ID: L12C2B8-06-TLN
% Mức: VDC
\begin{ex}
% ID: L12C2B8-06-TLN
% Mức: VDC
Trong dạng “Chuyển động Brown và bằng chứng về chuyển động phân tử”, Khí chiếm 5 lít, giãn nở thành 15 lít. Mật độ phân tử giảm bao nhiêu lần?
\shortans{3}
\loigiai{
$n_1/n_2=V_2/V_1=15/5=3$.
}
\end{ex}

% ===== Câu 24 =====
% ID: L12C2B8-06-TN
% Mức: VD
\begin{ex}
% ID: L12C2B8-06-TN
% Mức: VD
Trong dạng “Chuyển động Brown và bằng chứng về chuyển động phân tử”, Khi nhiệt độ khí tăng, chuyển động nhiệt của phân tử nhìn chung
\choice
{\True nhanh hơn}
{chậm hơn}
{không đổi}
{dừng lại}
\loigiai{
Đáp án A. Nhiệt độ tăng làm động năng tịnh tiến trung bình tăng.
}
\end{ex}

% ===== Câu 25 =====
% ID: L12C2B8-07-DS
% Mức: TH
\dangbt{Các luận điểm cơ bản của mô hình động học phân tử chất khí}
\begin{ex}
% ID: L12C2B8-07-DS
% Mức: TH
Xét các phát biểu sau trong dạng “Các luận điểm cơ bản của mô hình động học phân tử chất khí”.
\choiceTF
{\True Lực tương tác giữa các phân tử khí lí tưởng được coi là — phát biểu đúng là: không đáng kể trừ khi va chạm.}
{Lực tương tác giữa các phân tử khí lí tưởng được coi là — phát biểu luôn rất lớn là đúng.}
{\True Chuyển động Brown là chuyển động — phát biểu đúng là: hỗn loạn của hạt nhỏ lơ lửng do va chạm phân tử.}
{Chuyển động Brown là chuyển động — phát biểu rơi tự do của hạt bụi là đúng.}
\loigiai{
\begin{itemchoice}
\item (Đúng) Lực tương tác giữa các phân tử khí lí tưởng được coi là — phát biểu đúng là: không đáng kể trừ khi va chạm. (Vì): Ngoài thời gian va chạm rất ngắn, tương tác phân tử được bỏ qua. b) Sai: nội dung không phù hợp. c) Đúng: Các phân tử môi trường va chạm không cân bằng làm hạt nhỏ chuyển động Brown. d) Sai: nội dung không phù hợp
\item (Sai) Lực tương tác giữa các phân tử khí lí tưởng được coi là — phát biểu luôn rất lớn là đúng.
\item (Đúng) Chuyển động Brown là chuyển động — phát biểu đúng là: hỗn loạn của hạt nhỏ lơ lửng do va chạm phân tử.
\item (Sai) Chuyển động Brown là chuyển động — phát biểu rơi tự do của hạt bụi là đúng.
\end{itemchoice}
}
\end{ex}

% ===== Câu 26 =====
% ID: L12C2B8-07-TL
% Mức: TH
\begin{ex}
% ID: L12C2B8-07-TL
% Mức: TH
Trong dạng “Các luận điểm cơ bản của mô hình động học phân tử chất khí”, Trình bày các luận điểm chính của mô hình động học phân tử chất khí.
\choiceTF

\loigiai{
Chất khí gồm vô số phân tử rất nhỏ, chuyển động hỗn loạn không ngừng; khoảng cách giữa chúng lớn, tương tác ngoài va chạm không đáng kể; va chạm được coi là đàn hồi.
}
\end{ex}

% ===== Câu 27 =====
% ID: L12C2B8-07-TLN
% Mức: TH
\begin{ex}
% ID: L12C2B8-07-TLN
% Mức: TH
Trong dạng “Các luận điểm cơ bản của mô hình động học phân tử chất khí”, Một bình chứa $6\times10^{23}$ phân tử trong thể tích $0,02\,m^3$. Tính mật độ phân tử theo đơn vị $10^{25}\,m^{-3}$.
\shortans{3}
\loigiai{
$n_0=N/V=6\times10^{23}/0,02=3\times10^{25}\,m^{-3}$.
}
\end{ex}

% ===== Câu 28 =====
% ID: L12C2B8-07-TN
% Mức: VD
\dangbt{Chuyển động Brown và bằng chứng về chuyển động phân tử}
\begin{ex}
% ID: L12C2B8-07-TN
% Mức: VD
Trong dạng “Chuyển động Brown và bằng chứng về chuyển động phân tử”, Ba thông số trạng thái cơ bản của một lượng khí là
\choice
{\True áp suất, thể tích, nhiệt độ}
{khối lượng, trọng lượng, nhiệt lượng}
{vận tốc, gia tốc, thời gian}
{lực, công, công suất}
\loigiai{
Đáp án A. Trạng thái lượng khí xác định bởi p, $V$, T.
}
\end{ex}

% ===== Câu 29 =====
% ID: L12C2B8-08-DS
% Mức: VD
\dangbt{Các luận điểm cơ bản của mô hình động học phân tử chất khí}
\begin{ex}
% ID: L12C2B8-08-DS
% Mức: VD
Xét các phát biểu sau trong dạng “Các luận điểm cơ bản của mô hình động học phân tử chất khí”.
\choiceTF
{\True Áp suất khí lên thành bình sinh ra chủ yếu do — phát biểu đúng là: phân tử va chạm thành bình.}
{Áp suất khí lên thành bình sinh ra chủ yếu do — phát biểu trọng lượng bình là đúng.}
{\True Khi nhiệt độ khí tăng, chuyển động nhiệt của phân tử nhìn chung — phát biểu đúng là: nhanh hơn.}
{Khi nhiệt độ khí tăng, chuyển động nhiệt của phân tử nhìn chung — phát biểu không đổi là đúng.}
\loigiai{
\begin{itemchoice}
\item (Đúng) Áp suất khí lên thành bình sinh ra chủ yếu do — phát biểu đúng là: phân tử va chạm thành bình. (Vì): Xung lượng truyền trong các va chạm phân tử với thành bình tạo áp suất. b) Sai: nội dung không phù hợp. c) Đúng: Nhiệt độ tăng làm động năng tịnh tiến trung bình tăng. d) Sai: nội dung không phù hợp
\item (Sai) Áp suất khí lên thành bình sinh ra chủ yếu do — phát biểu trọng lượng bình là đúng.
\item (Đúng) Khi nhiệt độ khí tăng, chuyển động nhiệt của phân tử nhìn chung — phát biểu đúng là: nhanh hơn.
\item (Sai) Khi nhiệt độ khí tăng, chuyển động nhiệt của phân tử nhìn chung — phát biểu không đổi là đúng.
\end{itemchoice}
}
\end{ex}

% ===== Câu 30 =====
% ID: L12C2B8-08-TL
% Mức: TH
\begin{ex}
% ID: L12C2B8-08-TL
% Mức: TH
Trong dạng “Các luận điểm cơ bản của mô hình động học phân tử chất khí”, Giải thích chuyển động Brown.
\choiceTF

\loigiai{
Phân tử môi trường chuyển động hỗn loạn va chạm không cân bằng từ mọi phía vào hạt nhỏ, làm hạt chuyển động ngẫu nhiên.
}
\end{ex}

% ===== Câu 31 =====
% ID: L12C2B8-08-TLN
% Mức: TH
\begin{ex}
% ID: L12C2B8-08-TLN
% Mức: TH
Trong dạng “Các luận điểm cơ bản của mô hình động học phân tử chất khí”, Một lượng khí có thể tích 4 lít được nén còn 2 lít. Mật độ phân tử tăng bao nhiêu lần?
\shortans{2}
\loigiai{
Số phân tử không đổi nên mật độ tỉ lệ nghịch thể tích: $n_2/n_1=V_1/V_2=2$.
}
\end{ex}

% ===== Câu 32 =====
% ID: L12C2B8-08-TN
% Mức: VD
\dangbt{Chuyển động Brown và bằng chứng về chuyển động phân tử}
\begin{ex}
% ID: L12C2B8-08-TN
% Mức: VD
Trong dạng “Chuyển động Brown và bằng chứng về chuyển động phân tử”, Va chạm giữa các phân tử khí lí tưởng được xem là
\choice
{\True đàn hồi}
{mềm hoàn toàn}
{không bảo toàn động lượng}
{không xảy ra}
\loigiai{
Đáp án A. Mô hình khí lí tưởng coi các va chạm là đàn hồi.
}
\end{ex}

% ===== Câu 33 =====
% ID: L12C2B8-09-DS
% Mức: VD
\dangbt{Các luận điểm cơ bản của mô hình động học phân tử chất khí}
\begin{ex}
% ID: L12C2B8-09-DS
% Mức: VD
Xét các phát biểu sau trong dạng “Các luận điểm cơ bản của mô hình động học phân tử chất khí”.
\choiceTF
{\True Ba thông số trạng thái cơ bản của một lượng khí là — phát biểu đúng là: áp suất, thể tích, nhiệt độ.}
{Ba thông số trạng thái cơ bản của một lượng khí là — phát biểu khối lượng, trọng lượng, nhiệt lượng là đúng.}
{\True Va chạm giữa các phân tử khí lí tưởng được xem là — phát biểu đúng là: đàn hồi.}
{Va chạm giữa các phân tử khí lí tưởng được xem là — phát biểu không bảo toàn động lượng là đúng.}
\loigiai{
\begin{itemchoice}
\item (Đúng) Ba thông số trạng thái cơ bản của một lượng khí là — phát biểu đúng là: áp suất, thể tích, nhiệt độ. (Vì): Trạng thái lượng khí xác định bởi p, $V$, T. b) Sai: nội dung không phù hợp. c) Đúng: Mô hình khí lí tưởng coi các va chạm là đàn hồi. d) Sai: nội dung không phù hợp
\item (Sai) Ba thông số trạng thái cơ bản của một lượng khí là — phát biểu khối lượng, trọng lượng, nhiệt lượng là đúng.
\item (Đúng) Va chạm giữa các phân tử khí lí tưởng được xem là — phát biểu đúng là: đàn hồi.
\item (Sai) Va chạm giữa các phân tử khí lí tưởng được xem là — phát biểu không bảo toàn động lượng là đúng.
\end{itemchoice}
}
\end{ex}

% ===== Câu 34 =====
% ID: L12C2B8-09-TL
% Mức: VD
\begin{ex}
% ID: L12C2B8-09-TL
% Mức: VD
Trong dạng “Các luận điểm cơ bản của mô hình động học phân tử chất khí”, Giải thích vì sao chất khí gây áp suất lên thành bình.
\choiceTF

\loigiai{
Mỗi va chạm làm phân tử đổi động lượng và truyền xung lượng cho thành; tổng hợp rất nhiều va chạm tạo lực trung bình trên một đơn vị diện tích, tức áp suất.
}
\end{ex}

% ===== Câu 35 =====
% ID: L12C2B8-09-TLN
% Mức: VD
\begin{ex}
% ID: L12C2B8-09-TLN
% Mức: VD
Trong dạng “Các luận điểm cơ bản của mô hình động học phân tử chất khí”, Nhiệt độ tăng từ 300 $K$ lên 600 K. Động năng tịnh tiến trung bình tăng bao nhiêu lần?
\shortans{2}
\loigiai{
Động năng trung bình tỉ lệ với nhiệt độ tuyệt đối nên tăng $600/300=2$ lần.
}
\end{ex}

% ===== Câu 36 =====
% ID: L12C2B8-09-TN
% Mức: VDC
\dangbt{Chuyển động Brown và bằng chứng về chuyển động phân tử}
\begin{ex}
% ID: L12C2B8-09-TN
% Mức: VDC
Trong dạng “Chuyển động Brown và bằng chứng về chuyển động phân tử”, Khi nén khí ở nhiệt độ không đổi, mật độ phân tử
\choice
{\True tăng}
{giảm}
{không đổi}
{bằng không}
\loigiai{
Đáp án A. Cùng số phân tử trong thể tích nhỏ hơn nên mật độ phân tử tăng.
}
\end{ex}

% ===== Câu 37 =====
% ID: L12C2B8-10-DS
% Mức: VDC
\dangbt{Các luận điểm cơ bản của mô hình động học phân tử chất khí}
\begin{ex}
% ID: L12C2B8-10-DS
% Mức: VDC
Xét các phát biểu sau trong dạng “Các luận điểm cơ bản của mô hình động học phân tử chất khí”.
\choiceTF
{\True Khi nén khí ở nhiệt độ không đổi, mật độ phân tử — phát biểu đúng là: tăng.}
{Khi nén khí ở nhiệt độ không đổi, mật độ phân tử — phát biểu giảm là đúng.}
{\True Hiện tượng khuếch tán chứng tỏ các phân tử — phát biểu đúng là: chuyển động không ngừng.}
{Hiện tượng khuếch tán chứng tỏ các phân tử — phát biểu đứng yên là đúng.}
\loigiai{
\begin{itemchoice}
\item (Đúng) Khi nén khí ở nhiệt độ không đổi, mật độ phân tử — phát biểu đúng là: tăng. (Vì): Cùng số phân tử trong thể tích nhỏ hơn nên mật độ phân tử tăng. b) Sai: nội dung không phù hợp. c) Đúng: Khuếch tán là hệ quả trực tiếp của chuyển động nhiệt hỗn loạn. d) Sai: nội dung không phù hợp
\item (Sai) Khi nén khí ở nhiệt độ không đổi, mật độ phân tử — phát biểu giảm là đúng.
\item (Đúng) Hiện tượng khuếch tán chứng tỏ các phân tử — phát biểu đúng là: chuyển động không ngừng.
\item (Sai) Hiện tượng khuếch tán chứng tỏ các phân tử — phát biểu đứng yên là đúng.
\end{itemchoice}
}
\end{ex}

% ===== Câu 38 =====
% ID: L12C2B8-10-TL
% Mức: VD
\begin{ex}
% ID: L12C2B8-10-TL
% Mức: VD
Trong dạng “Các luận điểm cơ bản của mô hình động học phân tử chất khí”, So sánh chuyển động phân tử khi nhiệt độ tăng.
\choiceTF

\loigiai{
Nhiệt độ tuyệt đối tăng làm động năng tịnh tiến trung bình và tốc độ đặc trưng của phân tử tăng.
}
\end{ex}

% ===== Câu 39 =====
% ID: L12C2B8-10-TLN
% Mức: VD
\begin{ex}
% ID: L12C2B8-10-TLN
% Mức: VD
Trong dạng “Các luận điểm cơ bản của mô hình động học phân tử chất khí”, Khí có $N=1,2\times10^{24}$ phân tử và mật độ $3\times10^{25}\,m^{-3}$. Tính thể tích theo lít.
\shortans{40}
\loigiai{
$V=N/n_0=0,04\,m^3=40$ lít.
}
\end{ex}

% ===== Câu 40 =====
% ID: L12C2B8-10-TN
% Mức: VDC
\dangbt{Chuyển động Brown và bằng chứng về chuyển động phân tử}
\begin{ex}
% ID: L12C2B8-10-TN
% Mức: VDC
Trong dạng “Chuyển động Brown và bằng chứng về chuyển động phân tử”, Hiện tượng khuếch tán chứng tỏ các phân tử
\choice
{\True chuyển động không ngừng}
{có kích thước vĩ mô}
{đứng yên}
{không tương tác khi va chạm}
\loigiai{
Đáp án A. Khuếch tán là hệ quả trực tiếp của chuyển động nhiệt hỗn loạn.
}
\end{ex}

% ===== Câu 41 =====
% ID: L12C2B8-11-DS
% Mức: TH
\dangbt{Các thông số trạng thái của một lượng khí}
\begin{ex}
% ID: L12C2B8-11-DS
% Mức: TH
Xét các phát biểu sau trong dạng “Các thông số trạng thái của một lượng khí”.
\choiceTF
{\True Theo mô hình động học phân tử, các phân tử khí — phát biểu đúng là: chuyển động hỗn loạn không ngừng.}
{Theo mô hình động học phân tử, các phân tử khí — phát biểu đứng yên tại vị trí cân bằng là đúng.}
{\True Kích thước phân tử so với khoảng cách trung bình giữa chúng trong khí lí tưởng được coi là — phát biểu đúng là: rất nhỏ.}
{Kích thước phân tử so với khoảng cách trung bình giữa chúng trong khí lí tưởng được coi là — phát biểu bằng nhau là đúng.}
\loigiai{
\begin{itemchoice}
\item (Đúng) Theo mô hình động học phân tử, các phân tử khí — phát biểu đúng là: chuyển động hỗn loạn không ngừng. (Vì): Phân tử khí luôn chuyển động nhiệt hỗn loạn không ngừng. b) Sai: nội dung không phù hợp. c) Đúng: Trong mô hình khí lí tưởng, kích thước phân tử được bỏ qua so với khoảng cách giữa chúng. d) Sai: nội dung không phù hợp
\item (Sai) Theo mô hình động học phân tử, các phân tử khí — phát biểu đứng yên tại vị trí cân bằng là đúng.
\item (Đúng) Kích thước phân tử so với khoảng cách trung bình giữa chúng trong khí lí tưởng được coi là — phát biểu đúng là: rất nhỏ.
\item (Sai) Kích thước phân tử so với khoảng cách trung bình giữa chúng trong khí lí tưởng được coi là — phát biểu bằng nhau là đúng.
\end{itemchoice}
}
\end{ex}

% ===== Câu 42 =====
% ID: L12C2B8-11-TL
% Mức: VD
\dangbt{Các luận điểm cơ bản của mô hình động học phân tử chất khí}
\begin{ex}
% ID: L12C2B8-11-TL
% Mức: VD
Trong dạng “Các luận điểm cơ bản của mô hình động học phân tử chất khí”, Giải thích hiện tượng mùi nước hoa lan khắp phòng.
\choiceTF

\loigiai{
Phân tử nước hoa bay hơi rồi chuyển động nhiệt hỗn loạn, va chạm và khuếch tán từ nơi nồng độ cao sang thấp.
}
\end{ex}

% ===== Câu 43 =====
% ID: L12C2B8-11-TLN
% Mức: VD
\begin{ex}
% ID: L12C2B8-11-TLN
% Mức: VD
Trong dạng “Các luận điểm cơ bản của mô hình động học phân tử chất khí”, Có 2 mol khí. Lấy $N_A=6,02\times10^{23}$. Số phân tử theo đơn vị $10^{23}$ là bao nhiêu?
\shortans{12,04}
\loigiai{
$N=2N_A=12,04\times10^{23}$.
}
\end{ex}

% ===== Câu 44 =====
% ID: L12C2B8-11-TN
% Mức: NB
\begin{ex}
% ID: L12C2B8-11-TN
% Mức: NB
Trong dạng “Các luận điểm cơ bản của mô hình động học phân tử chất khí”, Theo mô hình động học phân tử, các phân tử khí
\choice
{\True chuyển động hỗn loạn không ngừng}
{đứng yên tại vị trí cân bằng}
{chỉ chuyển động khi được nung nóng}
{chuyển động cùng một hướng}
\loigiai{
Đáp án A. Phân tử khí luôn chuyển động nhiệt hỗn loạn không ngừng.
}
\end{ex}

% ===== Câu 45 =====
% ID: L12C2B8-12-DS
% Mức: TH
\dangbt{Các thông số trạng thái của một lượng khí}
\begin{ex}
% ID: L12C2B8-12-DS
% Mức: TH
Xét các phát biểu sau trong dạng “Các thông số trạng thái của một lượng khí”.
\choiceTF
{\True Lực tương tác giữa các phân tử khí lí tưởng được coi là — phát biểu đúng là: không đáng kể trừ khi va chạm.}
{Lực tương tác giữa các phân tử khí lí tưởng được coi là — phát biểu luôn rất lớn là đúng.}
{\True Chuyển động Brown là chuyển động — phát biểu đúng là: hỗn loạn của hạt nhỏ lơ lửng do va chạm phân tử.}
{Chuyển động Brown là chuyển động — phát biểu rơi tự do của hạt bụi là đúng.}
\loigiai{
\begin{itemchoice}
\item (Đúng) Lực tương tác giữa các phân tử khí lí tưởng được coi là — phát biểu đúng là: không đáng kể trừ khi va chạm. (Vì): Ngoài thời gian va chạm rất ngắn, tương tác phân tử được bỏ qua. b) Sai: nội dung không phù hợp. c) Đúng: Các phân tử môi trường va chạm không cân bằng làm hạt nhỏ chuyển động Brown. d) Sai: nội dung không phù hợp
\item (Sai) Lực tương tác giữa các phân tử khí lí tưởng được coi là — phát biểu luôn rất lớn là đúng.
\item (Đúng) Chuyển động Brown là chuyển động — phát biểu đúng là: hỗn loạn của hạt nhỏ lơ lửng do va chạm phân tử.
\item (Sai) Chuyển động Brown là chuyển động — phát biểu rơi tự do của hạt bụi là đúng.
\end{itemchoice}
}
\end{ex}

% ===== Câu 46 =====
% ID: L12C2B8-12-TL
% Mức: VDC
\dangbt{Các luận điểm cơ bản của mô hình động học phân tử chất khí}
\begin{ex}
% ID: L12C2B8-12-TL
% Mức: VDC
Trong dạng “Các luận điểm cơ bản của mô hình động học phân tử chất khí”, Nêu giới hạn của mô hình khí lí tưởng.
\choiceTF

\loigiai{
Mô hình phù hợp khi khí loãng, áp suất không quá cao và nhiệt độ không gần vùng hóa lỏng; khi đó thể tích riêng và tương tác phân tử có thể bỏ qua.
}
\end{ex}

% ===== Câu 47 =====
% ID: L12C2B8-12-TLN
% Mức: VDC
\begin{ex}
% ID: L12C2B8-12-TLN
% Mức: VDC
Trong dạng “Các luận điểm cơ bản của mô hình động học phân tử chất khí”, Khí chiếm 5 lít, giãn nở thành 15 lít. Mật độ phân tử giảm bao nhiêu lần?
\shortans{3}
\loigiai{
$n_1/n_2=V_2/V_1=15/5=3$.
}
\end{ex}

% ===== Câu 48 =====
% ID: L12C2B8-12-TN
% Mức: NB
\begin{ex}
% ID: L12C2B8-12-TN
% Mức: NB
Trong dạng “Các luận điểm cơ bản của mô hình động học phân tử chất khí”, Kích thước phân tử so với khoảng cách trung bình giữa chúng trong khí lí tưởng được coi là
\choice
{\True rất nhỏ}
{rất lớn}
{bằng nhau}
{không xác định}
\loigiai{
Đáp án A. Trong mô hình khí lí tưởng, kích thước phân tử được bỏ qua so với khoảng cách giữa chúng.
}
\end{ex}

% ===== Câu 49 =====
% ID: L12C2B8-13-DS
% Mức: VD
\dangbt{Các thông số trạng thái của một lượng khí}
\begin{ex}
% ID: L12C2B8-13-DS
% Mức: VD
Xét các phát biểu sau trong dạng “Các thông số trạng thái của một lượng khí”.
\choiceTF
{\True Áp suất khí lên thành bình sinh ra chủ yếu do — phát biểu đúng là: phân tử va chạm thành bình.}
{Áp suất khí lên thành bình sinh ra chủ yếu do — phát biểu trọng lượng bình là đúng.}
{\True Khi nhiệt độ khí tăng, chuyển động nhiệt của phân tử nhìn chung — phát biểu đúng là: nhanh hơn.}
{Khi nhiệt độ khí tăng, chuyển động nhiệt của phân tử nhìn chung — phát biểu không đổi là đúng.}
\loigiai{
\begin{itemchoice}
\item (Đúng) Áp suất khí lên thành bình sinh ra chủ yếu do — phát biểu đúng là: phân tử va chạm thành bình. (Vì): Xung lượng truyền trong các va chạm phân tử với thành bình tạo áp suất. b) Sai: nội dung không phù hợp. c) Đúng: Nhiệt độ tăng làm động năng tịnh tiến trung bình tăng. d) Sai: nội dung không phù hợp
\item (Sai) Áp suất khí lên thành bình sinh ra chủ yếu do — phát biểu trọng lượng bình là đúng.
\item (Đúng) Khi nhiệt độ khí tăng, chuyển động nhiệt của phân tử nhìn chung — phát biểu đúng là: nhanh hơn.
\item (Sai) Khi nhiệt độ khí tăng, chuyển động nhiệt của phân tử nhìn chung — phát biểu không đổi là đúng.
\end{itemchoice}
}
\end{ex}

% ===== Câu 50 =====
% ID: L12C2B8-13-TL
% Mức: TH
\begin{ex}
% ID: L12C2B8-13-TL
% Mức: TH
Trong dạng “Các thông số trạng thái của một lượng khí”, Trình bày các luận điểm chính của mô hình động học phân tử chất khí.
\choiceTF

\loigiai{
Chất khí gồm vô số phân tử rất nhỏ, chuyển động hỗn loạn không ngừng; khoảng cách giữa chúng lớn, tương tác ngoài va chạm không đáng kể; va chạm được coi là đàn hồi.
}
\end{ex}

% ===== Câu 51 =====
% ID: L12C2B8-13-TLN
% Mức: TH
\begin{ex}
% ID: L12C2B8-13-TLN
% Mức: TH
Trong dạng “Các thông số trạng thái của một lượng khí”, Một bình chứa $6\times10^{23}$ phân tử trong thể tích $0,02\,m^3$. Tính mật độ phân tử theo đơn vị $10^{25}\,m^{-3}$.
\shortans{3}
\loigiai{
$n_0=N/V=6\times10^{23}/0,02=3\times10^{25}\,m^{-3}$.
}
\end{ex}

% ===== Câu 52 =====
% ID: L12C2B8-13-TN
% Mức: TH
\dangbt{Các luận điểm cơ bản của mô hình động học phân tử chất khí}
\begin{ex}
% ID: L12C2B8-13-TN
% Mức: TH
Trong dạng “Các luận điểm cơ bản của mô hình động học phân tử chất khí”, Lực tương tác giữa các phân tử khí lí tưởng được coi là
\choice
{\True không đáng kể trừ khi va chạm}
{luôn rất lớn}
{lực hút không đổi}
{bằng trọng lực}
\loigiai{
Đáp án A. Ngoài thời gian va chạm rất ngắn, tương tác phân tử được bỏ qua.
}
\end{ex}

% ===== Câu 53 =====
% ID: L12C2B8-14-DS
% Mức: VD
\dangbt{Các thông số trạng thái của một lượng khí}
\begin{ex}
% ID: L12C2B8-14-DS
% Mức: VD
Xét các phát biểu sau trong dạng “Các thông số trạng thái của một lượng khí”.
\choiceTF
{\True Ba thông số trạng thái cơ bản của một lượng khí là — phát biểu đúng là: áp suất, thể tích, nhiệt độ.}
{Ba thông số trạng thái cơ bản của một lượng khí là — phát biểu khối lượng, trọng lượng, nhiệt lượng là đúng.}
{\True Va chạm giữa các phân tử khí lí tưởng được xem là — phát biểu đúng là: đàn hồi.}
{Va chạm giữa các phân tử khí lí tưởng được xem là — phát biểu không bảo toàn động lượng là đúng.}
\loigiai{
\begin{itemchoice}
\item (Đúng) Ba thông số trạng thái cơ bản của một lượng khí là — phát biểu đúng là: áp suất, thể tích, nhiệt độ. (Vì): Trạng thái lượng khí xác định bởi p, $V$, T. b) Sai: nội dung không phù hợp. c) Đúng: Mô hình khí lí tưởng coi các va chạm là đàn hồi. d) Sai: nội dung không phù hợp
\item (Sai) Ba thông số trạng thái cơ bản của một lượng khí là — phát biểu khối lượng, trọng lượng, nhiệt lượng là đúng.
\item (Đúng) Va chạm giữa các phân tử khí lí tưởng được xem là — phát biểu đúng là: đàn hồi.
\item (Sai) Va chạm giữa các phân tử khí lí tưởng được xem là — phát biểu không bảo toàn động lượng là đúng.
\end{itemchoice}
}
\end{ex}

% ===== Câu 54 =====
% ID: L12C2B8-14-TL
% Mức: TH
\begin{ex}
% ID: L12C2B8-14-TL
% Mức: TH
Trong dạng “Các thông số trạng thái của một lượng khí”, Giải thích chuyển động Brown.
\choiceTF

\loigiai{
Phân tử môi trường chuyển động hỗn loạn va chạm không cân bằng từ mọi phía vào hạt nhỏ, làm hạt chuyển động ngẫu nhiên.
}
\end{ex}

% ===== Câu 55 =====
% ID: L12C2B8-14-TLN
% Mức: TH
\begin{ex}
% ID: L12C2B8-14-TLN
% Mức: TH
Trong dạng “Các thông số trạng thái của một lượng khí”, Một lượng khí có thể tích 4 lít được nén còn 2 lít. Mật độ phân tử tăng bao nhiêu lần?
\shortans{2}
\loigiai{
Số phân tử không đổi nên mật độ tỉ lệ nghịch thể tích: $n_2/n_1=V_1/V_2=2$.
}
\end{ex}

% ===== Câu 56 =====
% ID: L12C2B8-14-TN
% Mức: TH
\dangbt{Các luận điểm cơ bản của mô hình động học phân tử chất khí}
\begin{ex}
% ID: L12C2B8-14-TN
% Mức: TH
Trong dạng “Các luận điểm cơ bản của mô hình động học phân tử chất khí”, Chuyển động Brown là chuyển động
\choice
{\True hỗn loạn của hạt nhỏ lơ lửng do va chạm phân tử}
{thẳng đều của phân tử khí}
{rơi tự do của hạt bụi}
{tuần hoàn của chất lỏng}
\loigiai{
Đáp án A. Các phân tử môi trường va chạm không cân bằng làm hạt nhỏ chuyển động Brown.
}
\end{ex}

% ===== Câu 57 =====
% ID: L12C2B8-15-DS
% Mức: VDC
\dangbt{Các thông số trạng thái của một lượng khí}
\begin{ex}
% ID: L12C2B8-15-DS
% Mức: VDC
Xét các phát biểu sau trong dạng “Các thông số trạng thái của một lượng khí”.
\choiceTF
{\True Khi nén khí ở nhiệt độ không đổi, mật độ phân tử — phát biểu đúng là: tăng.}
{Khi nén khí ở nhiệt độ không đổi, mật độ phân tử — phát biểu giảm là đúng.}
{\True Hiện tượng khuếch tán chứng tỏ các phân tử — phát biểu đúng là: chuyển động không ngừng.}
{Hiện tượng khuếch tán chứng tỏ các phân tử — phát biểu đứng yên là đúng.}
\loigiai{
\begin{itemchoice}
\item (Đúng) Khi nén khí ở nhiệt độ không đổi, mật độ phân tử — phát biểu đúng là: tăng. (Vì): Cùng số phân tử trong thể tích nhỏ hơn nên mật độ phân tử tăng. b) Sai: nội dung không phù hợp. c) Đúng: Khuếch tán là hệ quả trực tiếp của chuyển động nhiệt hỗn loạn. d) Sai: nội dung không phù hợp
\item (Sai) Khi nén khí ở nhiệt độ không đổi, mật độ phân tử — phát biểu giảm là đúng.
\item (Đúng) Hiện tượng khuếch tán chứng tỏ các phân tử — phát biểu đúng là: chuyển động không ngừng.
\item (Sai) Hiện tượng khuếch tán chứng tỏ các phân tử — phát biểu đứng yên là đúng.
\end{itemchoice}
}
\end{ex}

% ===== Câu 58 =====
% ID: L12C2B8-15-TL
% Mức: VD
\begin{ex}
% ID: L12C2B8-15-TL
% Mức: VD
Trong dạng “Các thông số trạng thái của một lượng khí”, Giải thích vì sao chất khí gây áp suất lên thành bình.
\choiceTF

\loigiai{
Mỗi va chạm làm phân tử đổi động lượng và truyền xung lượng cho thành; tổng hợp rất nhiều va chạm tạo lực trung bình trên một đơn vị diện tích, tức áp suất.
}
\end{ex}

% ===== Câu 59 =====
% ID: L12C2B8-15-TLN
% Mức: VD
\begin{ex}
% ID: L12C2B8-15-TLN
% Mức: VD
Trong dạng “Các thông số trạng thái của một lượng khí”, Nhiệt độ tăng từ 300 $K$ lên 600 K. Động năng tịnh tiến trung bình tăng bao nhiêu lần?
\shortans{2}
\loigiai{
Động năng trung bình tỉ lệ với nhiệt độ tuyệt đối nên tăng $600/300=2$ lần.
}
\end{ex}

% ===== Câu 60 =====
% ID: L12C2B8-15-TN
% Mức: TH
\dangbt{Các luận điểm cơ bản của mô hình động học phân tử chất khí}
\begin{ex}
% ID: L12C2B8-15-TN
% Mức: TH
Trong dạng “Các luận điểm cơ bản của mô hình động học phân tử chất khí”, Áp suất khí lên thành bình sinh ra chủ yếu do
\choice
{\True phân tử va chạm thành bình}
{trọng lượng bình}
{lực hút của Trái Đất}
{thành bình nóng lên}
\loigiai{
Đáp án A. Xung lượng truyền trong các va chạm phân tử với thành bình tạo áp suất.
}
\end{ex}

% ===== Câu 61 =====
% ID: L12C2B8-16-DS
% Mức: TH
\dangbt{Giải thích áp suất khí bằng va chạm phân tử}
\begin{ex}
% ID: L12C2B8-16-DS
% Mức: TH
Xét các phát biểu sau trong dạng “Giải thích áp suất khí bằng va chạm phân tử”.
\choiceTF
{\True Theo mô hình động học phân tử, các phân tử khí — phát biểu đúng là: chuyển động hỗn loạn không ngừng.}
{Theo mô hình động học phân tử, các phân tử khí — phát biểu đứng yên tại vị trí cân bằng là đúng.}
{\True Kích thước phân tử so với khoảng cách trung bình giữa chúng trong khí lí tưởng được coi là — phát biểu đúng là: rất nhỏ.}
{Kích thước phân tử so với khoảng cách trung bình giữa chúng trong khí lí tưởng được coi là — phát biểu bằng nhau là đúng.}
\loigiai{
\begin{itemchoice}
\item (Đúng) Theo mô hình động học phân tử, các phân tử khí — phát biểu đúng là: chuyển động hỗn loạn không ngừng. (Vì): Phân tử khí luôn chuyển động nhiệt hỗn loạn không ngừng. b) Sai: nội dung không phù hợp. c) Đúng: Trong mô hình khí lí tưởng, kích thước phân tử được bỏ qua so với khoảng cách giữa chúng. d) Sai: nội dung không phù hợp
\item (Sai) Theo mô hình động học phân tử, các phân tử khí — phát biểu đứng yên tại vị trí cân bằng là đúng.
\item (Đúng) Kích thước phân tử so với khoảng cách trung bình giữa chúng trong khí lí tưởng được coi là — phát biểu đúng là: rất nhỏ.
\item (Sai) Kích thước phân tử so với khoảng cách trung bình giữa chúng trong khí lí tưởng được coi là — phát biểu bằng nhau là đúng.
\end{itemchoice}
}
\end{ex}

% ===== Câu 62 =====
% ID: L12C2B8-16-TL
% Mức: VD
\dangbt{Các thông số trạng thái của một lượng khí}
\begin{ex}
% ID: L12C2B8-16-TL
% Mức: VD
Trong dạng “Các thông số trạng thái của một lượng khí”, So sánh chuyển động phân tử khi nhiệt độ tăng.
\choiceTF

\loigiai{
Nhiệt độ tuyệt đối tăng làm động năng tịnh tiến trung bình và tốc độ đặc trưng của phân tử tăng.
}
\end{ex}

% ===== Câu 63 =====
% ID: L12C2B8-16-TLN
% Mức: VD
\begin{ex}
% ID: L12C2B8-16-TLN
% Mức: VD
Trong dạng “Các thông số trạng thái của một lượng khí”, Khí có $N=1,2\times10^{24}$ phân tử và mật độ $3\times10^{25}\,m^{-3}$. Tính thể tích theo lít.
\shortans{40}
\loigiai{
$V=N/n_0=0,04\,m^3=40$ lít.
}
\end{ex}

% ===== Câu 64 =====
% ID: L12C2B8-16-TN
% Mức: VD
\dangbt{Các luận điểm cơ bản của mô hình động học phân tử chất khí}
\begin{ex}
% ID: L12C2B8-16-TN
% Mức: VD
Trong dạng “Các luận điểm cơ bản của mô hình động học phân tử chất khí”, Khi nhiệt độ khí tăng, chuyển động nhiệt của phân tử nhìn chung
\choice
{\True nhanh hơn}
{chậm hơn}
{không đổi}
{dừng lại}
\loigiai{
Đáp án A. Nhiệt độ tăng làm động năng tịnh tiến trung bình tăng.
}
\end{ex}

% ===== Câu 65 =====
% ID: L12C2B8-17-DS
% Mức: TH
\dangbt{Giải thích áp suất khí bằng va chạm phân tử}
\begin{ex}
% ID: L12C2B8-17-DS
% Mức: TH
Xét các phát biểu sau trong dạng “Giải thích áp suất khí bằng va chạm phân tử”.
\choiceTF
{\True Lực tương tác giữa các phân tử khí lí tưởng được coi là — phát biểu đúng là: không đáng kể trừ khi va chạm.}
{Lực tương tác giữa các phân tử khí lí tưởng được coi là — phát biểu luôn rất lớn là đúng.}
{\True Chuyển động Brown là chuyển động — phát biểu đúng là: hỗn loạn của hạt nhỏ lơ lửng do va chạm phân tử.}
{Chuyển động Brown là chuyển động — phát biểu rơi tự do của hạt bụi là đúng.}
\loigiai{
\begin{itemchoice}
\item (Đúng) Lực tương tác giữa các phân tử khí lí tưởng được coi là — phát biểu đúng là: không đáng kể trừ khi va chạm. (Vì): Ngoài thời gian va chạm rất ngắn, tương tác phân tử được bỏ qua. b) Sai: nội dung không phù hợp. c) Đúng: Các phân tử môi trường va chạm không cân bằng làm hạt nhỏ chuyển động Brown. d) Sai: nội dung không phù hợp
\item (Sai) Lực tương tác giữa các phân tử khí lí tưởng được coi là — phát biểu luôn rất lớn là đúng.
\item (Đúng) Chuyển động Brown là chuyển động — phát biểu đúng là: hỗn loạn của hạt nhỏ lơ lửng do va chạm phân tử.
\item (Sai) Chuyển động Brown là chuyển động — phát biểu rơi tự do của hạt bụi là đúng.
\end{itemchoice}
}
\end{ex}

% ===== Câu 66 =====
% ID: L12C2B8-17-TL
% Mức: VD
\dangbt{Các thông số trạng thái của một lượng khí}
\begin{ex}
% ID: L12C2B8-17-TL
% Mức: VD
Trong dạng “Các thông số trạng thái của một lượng khí”, Giải thích hiện tượng mùi nước hoa lan khắp phòng.
\choiceTF

\loigiai{
Phân tử nước hoa bay hơi rồi chuyển động nhiệt hỗn loạn, va chạm và khuếch tán từ nơi nồng độ cao sang thấp.
}
\end{ex}

% ===== Câu 67 =====
% ID: L12C2B8-17-TLN
% Mức: VD
\begin{ex}
% ID: L12C2B8-17-TLN
% Mức: VD
Trong dạng “Các thông số trạng thái của một lượng khí”, Có 2 mol khí. Lấy $N_A=6,02\times10^{23}$. Số phân tử theo đơn vị $10^{23}$ là bao nhiêu?
\shortans{12,04}
\loigiai{
$N=2N_A=12,04\times10^{23}$.
}
\end{ex}

% ===== Câu 68 =====
% ID: L12C2B8-17-TN
% Mức: VD
\dangbt{Các luận điểm cơ bản của mô hình động học phân tử chất khí}
\begin{ex}
% ID: L12C2B8-17-TN
% Mức: VD
Trong dạng “Các luận điểm cơ bản của mô hình động học phân tử chất khí”, Ba thông số trạng thái cơ bản của một lượng khí là
\choice
{\True áp suất, thể tích, nhiệt độ}
{khối lượng, trọng lượng, nhiệt lượng}
{vận tốc, gia tốc, thời gian}
{lực, công, công suất}
\loigiai{
Đáp án A. Trạng thái lượng khí xác định bởi p, $V$, T.
}
\end{ex}

% ===== Câu 69 =====
% ID: L12C2B8-18-DS
% Mức: VD
\dangbt{Giải thích áp suất khí bằng va chạm phân tử}
\begin{ex}
% ID: L12C2B8-18-DS
% Mức: VD
Xét các phát biểu sau trong dạng “Giải thích áp suất khí bằng va chạm phân tử”.
\choiceTF
{\True Áp suất khí lên thành bình sinh ra chủ yếu do — phát biểu đúng là: phân tử va chạm thành bình.}
{Áp suất khí lên thành bình sinh ra chủ yếu do — phát biểu trọng lượng bình là đúng.}
{\True Khi nhiệt độ khí tăng, chuyển động nhiệt của phân tử nhìn chung — phát biểu đúng là: nhanh hơn.}
{Khi nhiệt độ khí tăng, chuyển động nhiệt của phân tử nhìn chung — phát biểu không đổi là đúng.}
\loigiai{
\begin{itemchoice}
\item (Đúng) Áp suất khí lên thành bình sinh ra chủ yếu do — phát biểu đúng là: phân tử va chạm thành bình. (Vì): Xung lượng truyền trong các va chạm phân tử với thành bình tạo áp suất. b) Sai: nội dung không phù hợp. c) Đúng: Nhiệt độ tăng làm động năng tịnh tiến trung bình tăng. d) Sai: nội dung không phù hợp
\item (Sai) Áp suất khí lên thành bình sinh ra chủ yếu do — phát biểu trọng lượng bình là đúng.
\item (Đúng) Khi nhiệt độ khí tăng, chuyển động nhiệt của phân tử nhìn chung — phát biểu đúng là: nhanh hơn.
\item (Sai) Khi nhiệt độ khí tăng, chuyển động nhiệt của phân tử nhìn chung — phát biểu không đổi là đúng.
\end{itemchoice}
}
\end{ex}

% ===== Câu 70 =====
% ID: L12C2B8-18-TL
% Mức: VDC
\dangbt{Các thông số trạng thái của một lượng khí}
\begin{ex}
% ID: L12C2B8-18-TL
% Mức: VDC
Trong dạng “Các thông số trạng thái của một lượng khí”, Nêu giới hạn của mô hình khí lí tưởng.
\choiceTF

\loigiai{
Mô hình phù hợp khi khí loãng, áp suất không quá cao và nhiệt độ không gần vùng hóa lỏng; khi đó thể tích riêng và tương tác phân tử có thể bỏ qua.
}
\end{ex}

% ===== Câu 71 =====
% ID: L12C2B8-18-TLN
% Mức: VDC
\begin{ex}
% ID: L12C2B8-18-TLN
% Mức: VDC
Trong dạng “Các thông số trạng thái của một lượng khí”, Khí chiếm 5 lít, giãn nở thành 15 lít. Mật độ phân tử giảm bao nhiêu lần?
\shortans{3}
\loigiai{
$n_1/n_2=V_2/V_1=15/5=3$.
}
\end{ex}

% ===== Câu 72 =====
% ID: L12C2B8-18-TN
% Mức: VD
\dangbt{Các luận điểm cơ bản của mô hình động học phân tử chất khí}
\begin{ex}
% ID: L12C2B8-18-TN
% Mức: VD
Trong dạng “Các luận điểm cơ bản của mô hình động học phân tử chất khí”, Va chạm giữa các phân tử khí lí tưởng được xem là
\choice
{\True đàn hồi}
{mềm hoàn toàn}
{không bảo toàn động lượng}
{không xảy ra}
\loigiai{
Đáp án A. Mô hình khí lí tưởng coi các va chạm là đàn hồi.
}
\end{ex}

% ===== Câu 73 =====
% ID: L12C2B8-19-DS
% Mức: VD
\dangbt{Giải thích áp suất khí bằng va chạm phân tử}
\begin{ex}
% ID: L12C2B8-19-DS
% Mức: VD
Xét các phát biểu sau trong dạng “Giải thích áp suất khí bằng va chạm phân tử”.
\choiceTF
{\True Ba thông số trạng thái cơ bản của một lượng khí là — phát biểu đúng là: áp suất, thể tích, nhiệt độ.}
{Ba thông số trạng thái cơ bản của một lượng khí là — phát biểu khối lượng, trọng lượng, nhiệt lượng là đúng.}
{\True Va chạm giữa các phân tử khí lí tưởng được xem là — phát biểu đúng là: đàn hồi.}
{Va chạm giữa các phân tử khí lí tưởng được xem là — phát biểu không bảo toàn động lượng là đúng.}
\loigiai{
\begin{itemchoice}
\item (Đúng) Ba thông số trạng thái cơ bản của một lượng khí là — phát biểu đúng là: áp suất, thể tích, nhiệt độ. (Vì): Trạng thái lượng khí xác định bởi p, $V$, T. b) Sai: nội dung không phù hợp. c) Đúng: Mô hình khí lí tưởng coi các va chạm là đàn hồi. d) Sai: nội dung không phù hợp
\item (Sai) Ba thông số trạng thái cơ bản của một lượng khí là — phát biểu khối lượng, trọng lượng, nhiệt lượng là đúng.
\item (Đúng) Va chạm giữa các phân tử khí lí tưởng được xem là — phát biểu đúng là: đàn hồi.
\item (Sai) Va chạm giữa các phân tử khí lí tưởng được xem là — phát biểu không bảo toàn động lượng là đúng.
\end{itemchoice}
}
\end{ex}

% ===== Câu 74 =====
% ID: L12C2B8-19-TL
% Mức: TH
\begin{ex}
% ID: L12C2B8-19-TL
% Mức: TH
Trong dạng “Giải thích áp suất khí bằng va chạm phân tử”, Trình bày các luận điểm chính của mô hình động học phân tử chất khí.
\choiceTF

\loigiai{
Chất khí gồm vô số phân tử rất nhỏ, chuyển động hỗn loạn không ngừng; khoảng cách giữa chúng lớn, tương tác ngoài va chạm không đáng kể; va chạm được coi là đàn hồi.
}
\end{ex}

% ===== Câu 75 =====
% ID: L12C2B8-19-TLN
% Mức: TH
\begin{ex}
% ID: L12C2B8-19-TLN
% Mức: TH
Trong dạng “Giải thích áp suất khí bằng va chạm phân tử”, Một bình chứa $6\times10^{23}$ phân tử trong thể tích $0,02\,m^3$. Tính mật độ phân tử theo đơn vị $10^{25}\,m^{-3}$.
\shortans{3}
\loigiai{
$n_0=N/V=6\times10^{23}/0,02=3\times10^{25}\,m^{-3}$.
}
\end{ex}

% ===== Câu 76 =====
% ID: L12C2B8-19-TN
% Mức: VDC
\dangbt{Các luận điểm cơ bản của mô hình động học phân tử chất khí}
\begin{ex}
% ID: L12C2B8-19-TN
% Mức: VDC
Trong dạng “Các luận điểm cơ bản của mô hình động học phân tử chất khí”, Khi nén khí ở nhiệt độ không đổi, mật độ phân tử
\choice
{\True tăng}
{giảm}
{không đổi}
{bằng không}
\loigiai{
Đáp án A. Cùng số phân tử trong thể tích nhỏ hơn nên mật độ phân tử tăng.
}
\end{ex}

% ===== Câu 77 =====
% ID: L12C2B8-20-DS
% Mức: VDC
\dangbt{Giải thích áp suất khí bằng va chạm phân tử}
\begin{ex}
% ID: L12C2B8-20-DS
% Mức: VDC
Xét các phát biểu sau trong dạng “Giải thích áp suất khí bằng va chạm phân tử”.
\choiceTF
{\True Khi nén khí ở nhiệt độ không đổi, mật độ phân tử — phát biểu đúng là: tăng.}
{Khi nén khí ở nhiệt độ không đổi, mật độ phân tử — phát biểu giảm là đúng.}
{\True Hiện tượng khuếch tán chứng tỏ các phân tử — phát biểu đúng là: chuyển động không ngừng.}
{Hiện tượng khuếch tán chứng tỏ các phân tử — phát biểu đứng yên là đúng.}
\loigiai{
\begin{itemchoice}
\item (Đúng) Khi nén khí ở nhiệt độ không đổi, mật độ phân tử — phát biểu đúng là: tăng. (Vì): Cùng số phân tử trong thể tích nhỏ hơn nên mật độ phân tử tăng. b) Sai: nội dung không phù hợp. c) Đúng: Khuếch tán là hệ quả trực tiếp của chuyển động nhiệt hỗn loạn. d) Sai: nội dung không phù hợp
\item (Sai) Khi nén khí ở nhiệt độ không đổi, mật độ phân tử — phát biểu giảm là đúng.
\item (Đúng) Hiện tượng khuếch tán chứng tỏ các phân tử — phát biểu đúng là: chuyển động không ngừng.
\item (Sai) Hiện tượng khuếch tán chứng tỏ các phân tử — phát biểu đứng yên là đúng.
\end{itemchoice}
}
\end{ex}

% ===== Câu 78 =====
% ID: L12C2B8-20-TL
% Mức: TH
\begin{ex}
% ID: L12C2B8-20-TL
% Mức: TH
Trong dạng “Giải thích áp suất khí bằng va chạm phân tử”, Giải thích chuyển động Brown.
\choiceTF

\loigiai{
Phân tử môi trường chuyển động hỗn loạn va chạm không cân bằng từ mọi phía vào hạt nhỏ, làm hạt chuyển động ngẫu nhiên.
}
\end{ex}

% ===== Câu 79 =====
% ID: L12C2B8-20-TLN
% Mức: TH
\begin{ex}
% ID: L12C2B8-20-TLN
% Mức: TH
Trong dạng “Giải thích áp suất khí bằng va chạm phân tử”, Một lượng khí có thể tích 4 lít được nén còn 2 lít. Mật độ phân tử tăng bao nhiêu lần?
\shortans{2}
\loigiai{
Số phân tử không đổi nên mật độ tỉ lệ nghịch thể tích: $n_2/n_1=V_1/V_2=2$.
}
\end{ex}

% ===== Câu 80 =====
% ID: L12C2B8-20-TN
% Mức: VDC
\dangbt{Các luận điểm cơ bản của mô hình động học phân tử chất khí}
\begin{ex}
% ID: L12C2B8-20-TN
% Mức: VDC
Trong dạng “Các luận điểm cơ bản của mô hình động học phân tử chất khí”, Hiện tượng khuếch tán chứng tỏ các phân tử
\choice
{\True chuyển động không ngừng}
{có kích thước vĩ mô}
{đứng yên}
{không tương tác khi va chạm}
\loigiai{
Đáp án A. Khuếch tán là hệ quả trực tiếp của chuyển động nhiệt hỗn loạn.
}
\end{ex}

% ===== Câu 81 =====
% ID: L12C2B8-21-DS
% Mức: TH
\dangbt{Vận dụng mô hình động học phân tử vào thực tế}
\begin{ex}
% ID: L12C2B8-21-DS
% Mức: TH
Xét các phát biểu sau trong dạng “Vận dụng mô hình động học phân tử vào thực tế”.
\choiceTF
{\True Theo mô hình động học phân tử, các phân tử khí — phát biểu đúng là: chuyển động hỗn loạn không ngừng.}
{Theo mô hình động học phân tử, các phân tử khí — phát biểu đứng yên tại vị trí cân bằng là đúng.}
{\True Kích thước phân tử so với khoảng cách trung bình giữa chúng trong khí lí tưởng được coi là — phát biểu đúng là: rất nhỏ.}
{Kích thước phân tử so với khoảng cách trung bình giữa chúng trong khí lí tưởng được coi là — phát biểu bằng nhau là đúng.}
\loigiai{
\begin{itemchoice}
\item (Đúng) Theo mô hình động học phân tử, các phân tử khí — phát biểu đúng là: chuyển động hỗn loạn không ngừng. (Vì): Phân tử khí luôn chuyển động nhiệt hỗn loạn không ngừng. b) Sai: nội dung không phù hợp. c) Đúng: Trong mô hình khí lí tưởng, kích thước phân tử được bỏ qua so với khoảng cách giữa chúng. d) Sai: nội dung không phù hợp
\item (Sai) Theo mô hình động học phân tử, các phân tử khí — phát biểu đứng yên tại vị trí cân bằng là đúng.
\item (Đúng) Kích thước phân tử so với khoảng cách trung bình giữa chúng trong khí lí tưởng được coi là — phát biểu đúng là: rất nhỏ.
\item (Sai) Kích thước phân tử so với khoảng cách trung bình giữa chúng trong khí lí tưởng được coi là — phát biểu bằng nhau là đúng.
\end{itemchoice}
}
\end{ex}

% ===== Câu 82 =====
% ID: L12C2B8-21-TL
% Mức: VD
\dangbt{Giải thích áp suất khí bằng va chạm phân tử}
\begin{ex}
% ID: L12C2B8-21-TL
% Mức: VD
Trong dạng “Giải thích áp suất khí bằng va chạm phân tử”, Giải thích vì sao chất khí gây áp suất lên thành bình.
\choiceTF

\loigiai{
Mỗi va chạm làm phân tử đổi động lượng và truyền xung lượng cho thành; tổng hợp rất nhiều va chạm tạo lực trung bình trên một đơn vị diện tích, tức áp suất.
}
\end{ex}

% ===== Câu 83 =====
% ID: L12C2B8-21-TLN
% Mức: VD
\begin{ex}
% ID: L12C2B8-21-TLN
% Mức: VD
Trong dạng “Giải thích áp suất khí bằng va chạm phân tử”, Nhiệt độ tăng từ 300 $K$ lên 600 K. Động năng tịnh tiến trung bình tăng bao nhiêu lần?
\shortans{2}
\loigiai{
Động năng trung bình tỉ lệ với nhiệt độ tuyệt đối nên tăng $600/300=2$ lần.
}
\end{ex}

% ===== Câu 84 =====
% ID: L12C2B8-21-TN
% Mức: NB
\dangbt{Các thông số trạng thái của một lượng khí}
\begin{ex}
% ID: L12C2B8-21-TN
% Mức: NB
Trong dạng “Các thông số trạng thái của một lượng khí”, Theo mô hình động học phân tử, các phân tử khí
\choice
{\True chuyển động hỗn loạn không ngừng}
{đứng yên tại vị trí cân bằng}
{chỉ chuyển động khi được nung nóng}
{chuyển động cùng một hướng}
\loigiai{
Đáp án A. Phân tử khí luôn chuyển động nhiệt hỗn loạn không ngừng.
}
\end{ex}

% ===== Câu 85 =====
% ID: L12C2B8-22-DS
% Mức: TH
\dangbt{Vận dụng mô hình động học phân tử vào thực tế}
\begin{ex}
% ID: L12C2B8-22-DS
% Mức: TH
Xét các phát biểu sau trong dạng “Vận dụng mô hình động học phân tử vào thực tế”.
\choiceTF
{\True Lực tương tác giữa các phân tử khí lí tưởng được coi là — phát biểu đúng là: không đáng kể trừ khi va chạm.}
{Lực tương tác giữa các phân tử khí lí tưởng được coi là — phát biểu luôn rất lớn là đúng.}
{\True Chuyển động Brown là chuyển động — phát biểu đúng là: hỗn loạn của hạt nhỏ lơ lửng do va chạm phân tử.}
{Chuyển động Brown là chuyển động — phát biểu rơi tự do của hạt bụi là đúng.}
\loigiai{
\begin{itemchoice}
\item (Đúng) Lực tương tác giữa các phân tử khí lí tưởng được coi là — phát biểu đúng là: không đáng kể trừ khi va chạm. (Vì): Ngoài thời gian va chạm rất ngắn, tương tác phân tử được bỏ qua. b) Sai: nội dung không phù hợp. c) Đúng: Các phân tử môi trường va chạm không cân bằng làm hạt nhỏ chuyển động Brown. d) Sai: nội dung không phù hợp
\item (Sai) Lực tương tác giữa các phân tử khí lí tưởng được coi là — phát biểu luôn rất lớn là đúng.
\item (Đúng) Chuyển động Brown là chuyển động — phát biểu đúng là: hỗn loạn của hạt nhỏ lơ lửng do va chạm phân tử.
\item (Sai) Chuyển động Brown là chuyển động — phát biểu rơi tự do của hạt bụi là đúng.
\end{itemchoice}
}
\end{ex}

% ===== Câu 86 =====
% ID: L12C2B8-22-TL
% Mức: VD
\dangbt{Giải thích áp suất khí bằng va chạm phân tử}
\begin{ex}
% ID: L12C2B8-22-TL
% Mức: VD
Trong dạng “Giải thích áp suất khí bằng va chạm phân tử”, So sánh chuyển động phân tử khi nhiệt độ tăng.
\choiceTF

\loigiai{
Nhiệt độ tuyệt đối tăng làm động năng tịnh tiến trung bình và tốc độ đặc trưng của phân tử tăng.
}
\end{ex}

% ===== Câu 87 =====
% ID: L12C2B8-22-TLN
% Mức: VD
\begin{ex}
% ID: L12C2B8-22-TLN
% Mức: VD
Trong dạng “Giải thích áp suất khí bằng va chạm phân tử”, Khí có $N=1,2\times10^{24}$ phân tử và mật độ $3\times10^{25}\,m^{-3}$. Tính thể tích theo lít.
\shortans{40}
\loigiai{
$V=N/n_0=0,04\,m^3=40$ lít.
}
\end{ex}

% ===== Câu 88 =====
% ID: L12C2B8-22-TN
% Mức: NB
\dangbt{Các thông số trạng thái của một lượng khí}
\begin{ex}
% ID: L12C2B8-22-TN
% Mức: NB
Trong dạng “Các thông số trạng thái của một lượng khí”, Kích thước phân tử so với khoảng cách trung bình giữa chúng trong khí lí tưởng được coi là
\choice
{\True rất nhỏ}
{rất lớn}
{bằng nhau}
{không xác định}
\loigiai{
Đáp án A. Trong mô hình khí lí tưởng, kích thước phân tử được bỏ qua so với khoảng cách giữa chúng.
}
\end{ex}

% ===== Câu 89 =====
% ID: L12C2B8-23-DS
% Mức: VD
\dangbt{Vận dụng mô hình động học phân tử vào thực tế}
\begin{ex}
% ID: L12C2B8-23-DS
% Mức: VD
Xét các phát biểu sau trong dạng “Vận dụng mô hình động học phân tử vào thực tế”.
\choiceTF
{\True Áp suất khí lên thành bình sinh ra chủ yếu do — phát biểu đúng là: phân tử va chạm thành bình.}
{Áp suất khí lên thành bình sinh ra chủ yếu do — phát biểu trọng lượng bình là đúng.}
{\True Khi nhiệt độ khí tăng, chuyển động nhiệt của phân tử nhìn chung — phát biểu đúng là: nhanh hơn.}
{Khi nhiệt độ khí tăng, chuyển động nhiệt của phân tử nhìn chung — phát biểu không đổi là đúng.}
\loigiai{
\begin{itemchoice}
\item (Đúng) Áp suất khí lên thành bình sinh ra chủ yếu do — phát biểu đúng là: phân tử va chạm thành bình. (Vì): Xung lượng truyền trong các va chạm phân tử với thành bình tạo áp suất. b) Sai: nội dung không phù hợp. c) Đúng: Nhiệt độ tăng làm động năng tịnh tiến trung bình tăng. d) Sai: nội dung không phù hợp
\item (Sai) Áp suất khí lên thành bình sinh ra chủ yếu do — phát biểu trọng lượng bình là đúng.
\item (Đúng) Khi nhiệt độ khí tăng, chuyển động nhiệt của phân tử nhìn chung — phát biểu đúng là: nhanh hơn.
\item (Sai) Khi nhiệt độ khí tăng, chuyển động nhiệt của phân tử nhìn chung — phát biểu không đổi là đúng.
\end{itemchoice}
}
\end{ex}

% ===== Câu 90 =====
% ID: L12C2B8-23-TL
% Mức: VD
\dangbt{Giải thích áp suất khí bằng va chạm phân tử}
\begin{ex}
% ID: L12C2B8-23-TL
% Mức: VD
Trong dạng “Giải thích áp suất khí bằng va chạm phân tử”, Giải thích hiện tượng mùi nước hoa lan khắp phòng.
\choiceTF

\loigiai{
Phân tử nước hoa bay hơi rồi chuyển động nhiệt hỗn loạn, va chạm và khuếch tán từ nơi nồng độ cao sang thấp.
}
\end{ex}

% ===== Câu 91 =====
% ID: L12C2B8-23-TLN
% Mức: VD
\begin{ex}
% ID: L12C2B8-23-TLN
% Mức: VD
Trong dạng “Giải thích áp suất khí bằng va chạm phân tử”, Có 2 mol khí. Lấy $N_A=6,02\times10^{23}$. Số phân tử theo đơn vị $10^{23}$ là bao nhiêu?
\shortans{12,04}
\loigiai{
$N=2N_A=12,04\times10^{23}$.
}
\end{ex}

% ===== Câu 92 =====
% ID: L12C2B8-23-TN
% Mức: TH
\dangbt{Các thông số trạng thái của một lượng khí}
\begin{ex}
% ID: L12C2B8-23-TN
% Mức: TH
Trong dạng “Các thông số trạng thái của một lượng khí”, Lực tương tác giữa các phân tử khí lí tưởng được coi là
\choice
{\True không đáng kể trừ khi va chạm}
{luôn rất lớn}
{lực hút không đổi}
{bằng trọng lực}
\loigiai{
Đáp án A. Ngoài thời gian va chạm rất ngắn, tương tác phân tử được bỏ qua.
}
\end{ex}

% ===== Câu 93 =====
% ID: L12C2B8-24-DS
% Mức: VD
\dangbt{Vận dụng mô hình động học phân tử vào thực tế}
\begin{ex}
% ID: L12C2B8-24-DS
% Mức: VD
Xét các phát biểu sau trong dạng “Vận dụng mô hình động học phân tử vào thực tế”.
\choiceTF
{\True Ba thông số trạng thái cơ bản của một lượng khí là — phát biểu đúng là: áp suất, thể tích, nhiệt độ.}
{Ba thông số trạng thái cơ bản của một lượng khí là — phát biểu khối lượng, trọng lượng, nhiệt lượng là đúng.}
{\True Va chạm giữa các phân tử khí lí tưởng được xem là — phát biểu đúng là: đàn hồi.}
{Va chạm giữa các phân tử khí lí tưởng được xem là — phát biểu không bảo toàn động lượng là đúng.}
\loigiai{
\begin{itemchoice}
\item (Đúng) Ba thông số trạng thái cơ bản của một lượng khí là — phát biểu đúng là: áp suất, thể tích, nhiệt độ. (Vì): Trạng thái lượng khí xác định bởi p, $V$, T. b) Sai: nội dung không phù hợp. c) Đúng: Mô hình khí lí tưởng coi các va chạm là đàn hồi. d) Sai: nội dung không phù hợp
\item (Sai) Ba thông số trạng thái cơ bản của một lượng khí là — phát biểu khối lượng, trọng lượng, nhiệt lượng là đúng.
\item (Đúng) Va chạm giữa các phân tử khí lí tưởng được xem là — phát biểu đúng là: đàn hồi.
\item (Sai) Va chạm giữa các phân tử khí lí tưởng được xem là — phát biểu không bảo toàn động lượng là đúng.
\end{itemchoice}
}
\end{ex}

% ===== Câu 94 =====
% ID: L12C2B8-24-TL
% Mức: VDC
\dangbt{Giải thích áp suất khí bằng va chạm phân tử}
\begin{ex}
% ID: L12C2B8-24-TL
% Mức: VDC
Trong dạng “Giải thích áp suất khí bằng va chạm phân tử”, Nêu giới hạn của mô hình khí lí tưởng.
\choiceTF

\loigiai{
Mô hình phù hợp khi khí loãng, áp suất không quá cao và nhiệt độ không gần vùng hóa lỏng; khi đó thể tích riêng và tương tác phân tử có thể bỏ qua.
}
\end{ex}

% ===== Câu 95 =====
% ID: L12C2B8-24-TLN
% Mức: VDC
\begin{ex}
% ID: L12C2B8-24-TLN
% Mức: VDC
Trong dạng “Giải thích áp suất khí bằng va chạm phân tử”, Khí chiếm 5 lít, giãn nở thành 15 lít. Mật độ phân tử giảm bao nhiêu lần?
\shortans{3}
\loigiai{
$n_1/n_2=V_2/V_1=15/5=3$.
}
\end{ex}

% ===== Câu 96 =====
% ID: L12C2B8-24-TN
% Mức: TH
\dangbt{Các thông số trạng thái của một lượng khí}
\begin{ex}
% ID: L12C2B8-24-TN
% Mức: TH
Trong dạng “Các thông số trạng thái của một lượng khí”, Chuyển động Brown là chuyển động
\choice
{\True hỗn loạn của hạt nhỏ lơ lửng do va chạm phân tử}
{thẳng đều của phân tử khí}
{rơi tự do của hạt bụi}
{tuần hoàn của chất lỏng}
\loigiai{
Đáp án A. Các phân tử môi trường va chạm không cân bằng làm hạt nhỏ chuyển động Brown.
}
\end{ex}

% ===== Câu 97 =====
% ID: L12C2B8-25-DS
% Mức: VDC
\dangbt{Vận dụng mô hình động học phân tử vào thực tế}
\begin{ex}
% ID: L12C2B8-25-DS
% Mức: VDC
Xét các phát biểu sau trong dạng “Vận dụng mô hình động học phân tử vào thực tế”.
\choiceTF
{\True Khi nén khí ở nhiệt độ không đổi, mật độ phân tử — phát biểu đúng là: tăng.}
{Khi nén khí ở nhiệt độ không đổi, mật độ phân tử — phát biểu giảm là đúng.}
{\True Hiện tượng khuếch tán chứng tỏ các phân tử — phát biểu đúng là: chuyển động không ngừng.}
{Hiện tượng khuếch tán chứng tỏ các phân tử — phát biểu đứng yên là đúng.}
\loigiai{
\begin{itemchoice}
\item (Đúng) Khi nén khí ở nhiệt độ không đổi, mật độ phân tử — phát biểu đúng là: tăng. (Vì): Cùng số phân tử trong thể tích nhỏ hơn nên mật độ phân tử tăng. b) Sai: nội dung không phù hợp. c) Đúng: Khuếch tán là hệ quả trực tiếp của chuyển động nhiệt hỗn loạn. d) Sai: nội dung không phù hợp
\item (Sai) Khi nén khí ở nhiệt độ không đổi, mật độ phân tử — phát biểu giảm là đúng.
\item (Đúng) Hiện tượng khuếch tán chứng tỏ các phân tử — phát biểu đúng là: chuyển động không ngừng.
\item (Sai) Hiện tượng khuếch tán chứng tỏ các phân tử — phát biểu đứng yên là đúng.
\end{itemchoice}
}
\end{ex}

% ===== Câu 98 =====
% ID: L12C2B8-25-TL
% Mức: TH
\begin{ex}
% ID: L12C2B8-25-TL
% Mức: TH
Trong dạng “Vận dụng mô hình động học phân tử vào thực tế”, Trình bày các luận điểm chính của mô hình động học phân tử chất khí.
\choiceTF

\loigiai{
Chất khí gồm vô số phân tử rất nhỏ, chuyển động hỗn loạn không ngừng; khoảng cách giữa chúng lớn, tương tác ngoài va chạm không đáng kể; va chạm được coi là đàn hồi.
}
\end{ex}

% ===== Câu 99 =====
% ID: L12C2B8-25-TLN
% Mức: TH
\begin{ex}
% ID: L12C2B8-25-TLN
% Mức: TH
Trong dạng “Vận dụng mô hình động học phân tử vào thực tế”, Một bình chứa $6\times10^{23}$ phân tử trong thể tích $0,02\,m^3$. Tính mật độ phân tử theo đơn vị $10^{25}\,m^{-3}$.
\shortans{3}
\loigiai{
$n_0=N/V=6\times10^{23}/0,02=3\times10^{25}\,m^{-3}$.
}
\end{ex}

% ===== Câu 100 =====
% ID: L12C2B8-25-TN
% Mức: TH
\dangbt{Các thông số trạng thái của một lượng khí}
\begin{ex}
% ID: L12C2B8-25-TN
% Mức: TH
Trong dạng “Các thông số trạng thái của một lượng khí”, Áp suất khí lên thành bình sinh ra chủ yếu do
\choice
{\True phân tử va chạm thành bình}
{trọng lượng bình}
{lực hút của Trái Đất}
{thành bình nóng lên}
\loigiai{
Đáp án A. Xung lượng truyền trong các va chạm phân tử với thành bình tạo áp suất.
}
\end{ex}

% ===== Câu 101 =====
% ID: L12C2B8-26-DS
% Mức: TH
\dangbt{Đặc điểm cấu tạo và chuyển động phân tử ở thể khí}
\begin{ex}
% ID: L12C2B8-26-DS
% Mức: TH
Xét các phát biểu sau trong dạng “Đặc điểm cấu tạo và chuyển động phân tử ở thể khí”.
\choiceTF
{\True Theo mô hình động học phân tử, các phân tử khí — phát biểu đúng là: chuyển động hỗn loạn không ngừng.}
{Theo mô hình động học phân tử, các phân tử khí — phát biểu đứng yên tại vị trí cân bằng là đúng.}
{\True Kích thước phân tử so với khoảng cách trung bình giữa chúng trong khí lí tưởng được coi là — phát biểu đúng là: rất nhỏ.}
{Kích thước phân tử so với khoảng cách trung bình giữa chúng trong khí lí tưởng được coi là — phát biểu bằng nhau là đúng.}
\loigiai{
\begin{itemchoice}
\item (Đúng) Theo mô hình động học phân tử, các phân tử khí — phát biểu đúng là: chuyển động hỗn loạn không ngừng. (Vì): Phân tử khí luôn chuyển động nhiệt hỗn loạn không ngừng. b) Sai: nội dung không phù hợp. c) Đúng: Trong mô hình khí lí tưởng, kích thước phân tử được bỏ qua so với khoảng cách giữa chúng. d) Sai: nội dung không phù hợp
\item (Sai) Theo mô hình động học phân tử, các phân tử khí — phát biểu đứng yên tại vị trí cân bằng là đúng.
\item (Đúng) Kích thước phân tử so với khoảng cách trung bình giữa chúng trong khí lí tưởng được coi là — phát biểu đúng là: rất nhỏ.
\item (Sai) Kích thước phân tử so với khoảng cách trung bình giữa chúng trong khí lí tưởng được coi là — phát biểu bằng nhau là đúng.
\end{itemchoice}
}
\end{ex}

% ===== Câu 102 =====
% ID: L12C2B8-26-TL
% Mức: TH
\dangbt{Vận dụng mô hình động học phân tử vào thực tế}
\begin{ex}
% ID: L12C2B8-26-TL
% Mức: TH
Trong dạng “Vận dụng mô hình động học phân tử vào thực tế”, Giải thích chuyển động Brown.
\choiceTF

\loigiai{
Phân tử môi trường chuyển động hỗn loạn va chạm không cân bằng từ mọi phía vào hạt nhỏ, làm hạt chuyển động ngẫu nhiên.
}
\end{ex}

% ===== Câu 103 =====
% ID: L12C2B8-26-TLN
% Mức: TH
\begin{ex}
% ID: L12C2B8-26-TLN
% Mức: TH
Trong dạng “Vận dụng mô hình động học phân tử vào thực tế”, Một lượng khí có thể tích 4 lít được nén còn 2 lít. Mật độ phân tử tăng bao nhiêu lần?
\shortans{2}
\loigiai{
Số phân tử không đổi nên mật độ tỉ lệ nghịch thể tích: $n_2/n_1=V_1/V_2=2$.
}
\end{ex}

% ===== Câu 104 =====
% ID: L12C2B8-26-TN
% Mức: VD
\dangbt{Các thông số trạng thái của một lượng khí}
\begin{ex}
% ID: L12C2B8-26-TN
% Mức: VD
Trong dạng “Các thông số trạng thái của một lượng khí”, Khi nhiệt độ khí tăng, chuyển động nhiệt của phân tử nhìn chung
\choice
{\True nhanh hơn}
{chậm hơn}
{không đổi}
{dừng lại}
\loigiai{
Đáp án A. Nhiệt độ tăng làm động năng tịnh tiến trung bình tăng.
}
\end{ex}

% ===== Câu 105 =====
% ID: L12C2B8-27-DS
% Mức: TH
\dangbt{Đặc điểm cấu tạo và chuyển động phân tử ở thể khí}
\begin{ex}
% ID: L12C2B8-27-DS
% Mức: TH
Xét các phát biểu sau trong dạng “Đặc điểm cấu tạo và chuyển động phân tử ở thể khí”.
\choiceTF
{\True Lực tương tác giữa các phân tử khí lí tưởng được coi là — phát biểu đúng là: không đáng kể trừ khi va chạm.}
{Lực tương tác giữa các phân tử khí lí tưởng được coi là — phát biểu luôn rất lớn là đúng.}
{\True Chuyển động Brown là chuyển động — phát biểu đúng là: hỗn loạn của hạt nhỏ lơ lửng do va chạm phân tử.}
{Chuyển động Brown là chuyển động — phát biểu rơi tự do của hạt bụi là đúng.}
\loigiai{
\begin{itemchoice}
\item (Đúng) Lực tương tác giữa các phân tử khí lí tưởng được coi là — phát biểu đúng là: không đáng kể trừ khi va chạm. (Vì): Ngoài thời gian va chạm rất ngắn, tương tác phân tử được bỏ qua. b) Sai: nội dung không phù hợp. c) Đúng: Các phân tử môi trường va chạm không cân bằng làm hạt nhỏ chuyển động Brown. d) Sai: nội dung không phù hợp
\item (Sai) Lực tương tác giữa các phân tử khí lí tưởng được coi là — phát biểu luôn rất lớn là đúng.
\item (Đúng) Chuyển động Brown là chuyển động — phát biểu đúng là: hỗn loạn của hạt nhỏ lơ lửng do va chạm phân tử.
\item (Sai) Chuyển động Brown là chuyển động — phát biểu rơi tự do của hạt bụi là đúng.
\end{itemchoice}
}
\end{ex}

% ===== Câu 106 =====
% ID: L12C2B8-27-TL
% Mức: VD
\dangbt{Vận dụng mô hình động học phân tử vào thực tế}
\begin{ex}
% ID: L12C2B8-27-TL
% Mức: VD
Trong dạng “Vận dụng mô hình động học phân tử vào thực tế”, Giải thích vì sao chất khí gây áp suất lên thành bình.
\choiceTF

\loigiai{
Mỗi va chạm làm phân tử đổi động lượng và truyền xung lượng cho thành; tổng hợp rất nhiều va chạm tạo lực trung bình trên một đơn vị diện tích, tức áp suất.
}
\end{ex}

% ===== Câu 107 =====
% ID: L12C2B8-27-TLN
% Mức: VD
\begin{ex}
% ID: L12C2B8-27-TLN
% Mức: VD
Trong dạng “Vận dụng mô hình động học phân tử vào thực tế”, Nhiệt độ tăng từ 300 $K$ lên 600 K. Động năng tịnh tiến trung bình tăng bao nhiêu lần?
\shortans{2}
\loigiai{
Động năng trung bình tỉ lệ với nhiệt độ tuyệt đối nên tăng $600/300=2$ lần.
}
\end{ex}

% ===== Câu 108 =====
% ID: L12C2B8-27-TN
% Mức: VD
\dangbt{Các thông số trạng thái của một lượng khí}
\begin{ex}
% ID: L12C2B8-27-TN
% Mức: VD
Trong dạng “Các thông số trạng thái của một lượng khí”, Ba thông số trạng thái cơ bản của một lượng khí là
\choice
{\True áp suất, thể tích, nhiệt độ}
{khối lượng, trọng lượng, nhiệt lượng}
{vận tốc, gia tốc, thời gian}
{lực, công, công suất}
\loigiai{
Đáp án A. Trạng thái lượng khí xác định bởi p, $V$, T.
}
\end{ex}

% ===== Câu 109 =====
% ID: L12C2B8-28-DS
% Mức: VD
\dangbt{Đặc điểm cấu tạo và chuyển động phân tử ở thể khí}
\begin{ex}
% ID: L12C2B8-28-DS
% Mức: VD
Xét các phát biểu sau trong dạng “Đặc điểm cấu tạo và chuyển động phân tử ở thể khí”.
\choiceTF
{\True Áp suất khí lên thành bình sinh ra chủ yếu do — phát biểu đúng là: phân tử va chạm thành bình.}
{Áp suất khí lên thành bình sinh ra chủ yếu do — phát biểu trọng lượng bình là đúng.}
{\True Khi nhiệt độ khí tăng, chuyển động nhiệt của phân tử nhìn chung — phát biểu đúng là: nhanh hơn.}
{Khi nhiệt độ khí tăng, chuyển động nhiệt của phân tử nhìn chung — phát biểu không đổi là đúng.}
\loigiai{
\begin{itemchoice}
\item (Đúng) Áp suất khí lên thành bình sinh ra chủ yếu do — phát biểu đúng là: phân tử va chạm thành bình. (Vì): Xung lượng truyền trong các va chạm phân tử với thành bình tạo áp suất. b) Sai: nội dung không phù hợp. c) Đúng: Nhiệt độ tăng làm động năng tịnh tiến trung bình tăng. d) Sai: nội dung không phù hợp
\item (Sai) Áp suất khí lên thành bình sinh ra chủ yếu do — phát biểu trọng lượng bình là đúng.
\item (Đúng) Khi nhiệt độ khí tăng, chuyển động nhiệt của phân tử nhìn chung — phát biểu đúng là: nhanh hơn.
\item (Sai) Khi nhiệt độ khí tăng, chuyển động nhiệt của phân tử nhìn chung — phát biểu không đổi là đúng.
\end{itemchoice}
}
\end{ex}

% ===== Câu 110 =====
% ID: L12C2B8-28-TL
% Mức: VD
\dangbt{Vận dụng mô hình động học phân tử vào thực tế}
\begin{ex}
% ID: L12C2B8-28-TL
% Mức: VD
Trong dạng “Vận dụng mô hình động học phân tử vào thực tế”, So sánh chuyển động phân tử khi nhiệt độ tăng.
\choiceTF

\loigiai{
Nhiệt độ tuyệt đối tăng làm động năng tịnh tiến trung bình và tốc độ đặc trưng của phân tử tăng.
}
\end{ex}

% ===== Câu 111 =====
% ID: L12C2B8-28-TLN
% Mức: VD
\begin{ex}
% ID: L12C2B8-28-TLN
% Mức: VD
Trong dạng “Vận dụng mô hình động học phân tử vào thực tế”, Khí có $N=1,2\times10^{24}$ phân tử và mật độ $3\times10^{25}\,m^{-3}$. Tính thể tích theo lít.
\shortans{40}
\loigiai{
$V=N/n_0=0,04\,m^3=40$ lít.
}
\end{ex}

% ===== Câu 112 =====
% ID: L12C2B8-28-TN
% Mức: VD
\dangbt{Các thông số trạng thái của một lượng khí}
\begin{ex}
% ID: L12C2B8-28-TN
% Mức: VD
Trong dạng “Các thông số trạng thái của một lượng khí”, Va chạm giữa các phân tử khí lí tưởng được xem là
\choice
{\True đàn hồi}
{mềm hoàn toàn}
{không bảo toàn động lượng}
{không xảy ra}
\loigiai{
Đáp án A. Mô hình khí lí tưởng coi các va chạm là đàn hồi.
}
\end{ex}

% ===== Câu 113 =====
% ID: L12C2B8-29-DS
% Mức: VD
\dangbt{Đặc điểm cấu tạo và chuyển động phân tử ở thể khí}
\begin{ex}
% ID: L12C2B8-29-DS
% Mức: VD
Xét các phát biểu sau trong dạng “Đặc điểm cấu tạo và chuyển động phân tử ở thể khí”.
\choiceTF
{\True Ba thông số trạng thái cơ bản của một lượng khí là — phát biểu đúng là: áp suất, thể tích, nhiệt độ.}
{Ba thông số trạng thái cơ bản của một lượng khí là — phát biểu khối lượng, trọng lượng, nhiệt lượng là đúng.}
{\True Va chạm giữa các phân tử khí lí tưởng được xem là — phát biểu đúng là: đàn hồi.}
{Va chạm giữa các phân tử khí lí tưởng được xem là — phát biểu không bảo toàn động lượng là đúng.}
\loigiai{
\begin{itemchoice}
\item (Đúng) Ba thông số trạng thái cơ bản của một lượng khí là — phát biểu đúng là: áp suất, thể tích, nhiệt độ. (Vì): Trạng thái lượng khí xác định bởi p, $V$, T. b) Sai: nội dung không phù hợp. c) Đúng: Mô hình khí lí tưởng coi các va chạm là đàn hồi. d) Sai: nội dung không phù hợp
\item (Sai) Ba thông số trạng thái cơ bản của một lượng khí là — phát biểu khối lượng, trọng lượng, nhiệt lượng là đúng.
\item (Đúng) Va chạm giữa các phân tử khí lí tưởng được xem là — phát biểu đúng là: đàn hồi.
\item (Sai) Va chạm giữa các phân tử khí lí tưởng được xem là — phát biểu không bảo toàn động lượng là đúng.
\end{itemchoice}
}
\end{ex}

% ===== Câu 114 =====
% ID: L12C2B8-29-TL
% Mức: VD
\dangbt{Vận dụng mô hình động học phân tử vào thực tế}
\begin{ex}
% ID: L12C2B8-29-TL
% Mức: VD
Trong dạng “Vận dụng mô hình động học phân tử vào thực tế”, Giải thích hiện tượng mùi nước hoa lan khắp phòng.
\choiceTF

\loigiai{
Phân tử nước hoa bay hơi rồi chuyển động nhiệt hỗn loạn, va chạm và khuếch tán từ nơi nồng độ cao sang thấp.
}
\end{ex}

% ===== Câu 115 =====
% ID: L12C2B8-29-TLN
% Mức: VD
\begin{ex}
% ID: L12C2B8-29-TLN
% Mức: VD
Trong dạng “Vận dụng mô hình động học phân tử vào thực tế”, Có 2 mol khí. Lấy $N_A=6,02\times10^{23}$. Số phân tử theo đơn vị $10^{23}$ là bao nhiêu?
\shortans{12,04}
\loigiai{
$N=2N_A=12,04\times10^{23}$.
}
\end{ex}

% ===== Câu 116 =====
% ID: L12C2B8-29-TN
% Mức: VDC
\dangbt{Các thông số trạng thái của một lượng khí}
\begin{ex}
% ID: L12C2B8-29-TN
% Mức: VDC
Trong dạng “Các thông số trạng thái của một lượng khí”, Khi nén khí ở nhiệt độ không đổi, mật độ phân tử
\choice
{\True tăng}
{giảm}
{không đổi}
{bằng không}
\loigiai{
Đáp án A. Cùng số phân tử trong thể tích nhỏ hơn nên mật độ phân tử tăng.
}
\end{ex}

% ===== Câu 117 =====
% ID: L12C2B8-30-DS
% Mức: VDC
\dangbt{Đặc điểm cấu tạo và chuyển động phân tử ở thể khí}
\begin{ex}
% ID: L12C2B8-30-DS
% Mức: VDC
Xét các phát biểu sau trong dạng “Đặc điểm cấu tạo và chuyển động phân tử ở thể khí”.
\choiceTF
{\True Khi nén khí ở nhiệt độ không đổi, mật độ phân tử — phát biểu đúng là: tăng.}
{Khi nén khí ở nhiệt độ không đổi, mật độ phân tử — phát biểu giảm là đúng.}
{\True Hiện tượng khuếch tán chứng tỏ các phân tử — phát biểu đúng là: chuyển động không ngừng.}
{Hiện tượng khuếch tán chứng tỏ các phân tử — phát biểu đứng yên là đúng.}
\loigiai{
\begin{itemchoice}
\item (Đúng) Khi nén khí ở nhiệt độ không đổi, mật độ phân tử — phát biểu đúng là: tăng. (Vì): Cùng số phân tử trong thể tích nhỏ hơn nên mật độ phân tử tăng. b) Sai: nội dung không phù hợp. c) Đúng: Khuếch tán là hệ quả trực tiếp của chuyển động nhiệt hỗn loạn. d) Sai: nội dung không phù hợp
\item (Sai) Khi nén khí ở nhiệt độ không đổi, mật độ phân tử — phát biểu giảm là đúng.
\item (Đúng) Hiện tượng khuếch tán chứng tỏ các phân tử — phát biểu đúng là: chuyển động không ngừng.
\item (Sai) Hiện tượng khuếch tán chứng tỏ các phân tử — phát biểu đứng yên là đúng.
\end{itemchoice}
}
\end{ex}

% ===== Câu 118 =====
% ID: L12C2B8-30-TL
% Mức: VDC
\dangbt{Vận dụng mô hình động học phân tử vào thực tế}
\begin{ex}
% ID: L12C2B8-30-TL
% Mức: VDC
Trong dạng “Vận dụng mô hình động học phân tử vào thực tế”, Nêu giới hạn của mô hình khí lí tưởng.
\choiceTF

\loigiai{
Mô hình phù hợp khi khí loãng, áp suất không quá cao và nhiệt độ không gần vùng hóa lỏng; khi đó thể tích riêng và tương tác phân tử có thể bỏ qua.
}
\end{ex}

% ===== Câu 119 =====
% ID: L12C2B8-30-TLN
% Mức: VDC
\begin{ex}
% ID: L12C2B8-30-TLN
% Mức: VDC
Trong dạng “Vận dụng mô hình động học phân tử vào thực tế”, Khí chiếm 5 lít, giãn nở thành 15 lít. Mật độ phân tử giảm bao nhiêu lần?
\shortans{3}
\loigiai{
$n_1/n_2=V_2/V_1=15/5=3$.
}
\end{ex}

% ===== Câu 120 =====
% ID: L12C2B8-30-TN
% Mức: VDC
\dangbt{Các thông số trạng thái của một lượng khí}
\begin{ex}
% ID: L12C2B8-30-TN
% Mức: VDC
Trong dạng “Các thông số trạng thái của một lượng khí”, Hiện tượng khuếch tán chứng tỏ các phân tử
\choice
{\True chuyển động không ngừng}
{có kích thước vĩ mô}
{đứng yên}
{không tương tác khi va chạm}
\loigiai{
Đáp án A. Khuếch tán là hệ quả trực tiếp của chuyển động nhiệt hỗn loạn.
}
\end{ex}

% ===== Câu 121 =====
% ID: L12C2B8-31-TL
% Mức: TH
\dangbt{Chuyển động Brown và bằng chứng về chuyển động phân tử}
\begin{ex}
% ID: L12C2B8-31-TL
% Mức: TH
Xét các phát biểu sau trong dạng “Chuyển động Brown và bằng chứng về chuyển động phân tử”.
\choiceTF
{\True Theo mô hình động học phân tử chất khí, các phân tử khí luôn thực hiện chuyển động hỗn loạn không ngừng.}
{Theo mô hình động học phân tử chất khí, khẳng định cho rằng các phân tử khí luôn nằm yên tại vị trí cân bằng của chúng là đúng.}
{\True Kích thước của mỗi phân tử khí so với khoảng cách trung bình giữa các phân tử trong khí lí tưởng là rất nhỏ.}
{Nhận định khẳng định rằng kích thước của từng phân tử trong khí lí tưởng bằng với khoảng cách trung bình giữa các phân tử ấy là đúng.}
\loigiai{
\begin{itemchoice}
\item (Đúng) Theo mô hình động học phân tử chất khí, các phân tử khí luôn thực hiện chuyển động hỗn loạn không ngừng. (Vì): Theo mô hình động học phân tử chất khí, các phân tử khí luôn thực hiện chuyển động nhiệt hỗn loạn không ngừng, đây là một trong những luận điểm cơ bản của thuyết động học phân tử
\item (Sai) Theo mô hình động học phân tử chất khí, khẳng định cho rằng các phân tử khí luôn nằm yên tại vị trí cân bằng của chúng là đúng. (Vì): Theo mô hình động học phân tử chất khí, các phân tử khí không nằm yên tại vị trí cân bằng mà luôn chuyển động hỗn loạn không ngừng. Khẳng định các phân tử khí nằm yên tại vị trí cân bằng là hoàn toàn sai
\item (Đúng) Kích thước của mỗi phân tử khí so với khoảng cách trung bình giữa các phân tử trong khí lí tưởng là rất nhỏ. (Vì): Trong mô hình khí lí tưởng, kích thước của mỗi phân tử khí được coi là rất nhỏ so với khoảng cách trung bình giữa các phân tử, do đó thể tích riêng của các phân tử được bỏ qua
\item (Sai) Nhận định khẳng định rằng kích thước của từng phân tử trong khí lí tưởng bằng với khoảng cách trung bình giữa các phân tử ấy là đúng. (Vì): Trong khí lí tưởng, kích thước của từng phân tử là rất nhỏ so với khoảng cách trung bình giữa các phân tử, chứ không bằng nhau. Khẳng định kích thước phân tử bằng khoảng cách trung bình giữa các phân tử là sai
\end{itemchoice}
}
\end{ex}

% ===== Câu 122 =====
% ID: L12C2B8-31-TLN
% Mức: TH
\dangbt{Đặc điểm cấu tạo và chuyển động phân tử ở thể khí}
\begin{ex}
% ID: L12C2B8-31-TLN
% Mức: TH
Trong dạng “Đặc điểm cấu tạo và chuyển động phân tử ở thể khí”, Một bình chứa $6\times10^{23}$ phân tử trong thể tích $0,02\,m^3$. Tính mật độ phân tử theo đơn vị $10^{25}\,m^{-3}$.
\shortans{3}
\loigiai{
$n_0=N/V=6\times10^{23}/0,02=3\times10^{25}\,m^{-3}$.
}
\end{ex}

% ===== Câu 123 =====
% ID: L12C2B8-31-TN
% Mức: NB
\dangbt{Giải thích áp suất khí bằng va chạm phân tử}
\begin{ex}
% ID: L12C2B8-31-TN
% Mức: NB
Trong dạng “Giải thích áp suất khí bằng va chạm phân tử”, Theo mô hình động học phân tử, các phân tử khí
\choice
{\True chuyển động hỗn loạn không ngừng}
{đứng yên tại vị trí cân bằng}
{chỉ chuyển động khi được nung nóng}
{chuyển động cùng một hướng}
\loigiai{
Số phân tử $N = 6 \times 10^{23}$, thể tích $V = 0,02 \, \mathrm{m^3}$.
Mật độ phân tử $n = \frac{N}{V} = \frac{6 \times 10^{23}}{0,02} = 3 \times 10^{25} \, \mathrm{m^{-3}}$.
Giá trị mật độ phân tử theo đơn vị $10^{25} \, \mathrm{m^{-3}}$ là 3.

Đáp số: 3
}
\end{ex}

% ===== Câu 124 =====
% ID: L12C2B8-32-TL
% Mức: TH
\dangbt{Đặc điểm cấu tạo và chuyển động phân tử ở thể khí}
\begin{ex}
% ID: L12C2B8-32-TL
% Mức: TH
Trong dạng “Đặc điểm cấu tạo và chuyển động phân tử ở thể khí”, Giải thích chuyển động Brown.
\choiceTF

\loigiai{
Phân tử môi trường chuyển động hỗn loạn va chạm không cân bằng từ mọi phía vào hạt nhỏ, làm hạt chuyển động ngẫu nhiên.
}
\end{ex}

% ===== Câu 125 =====
% ID: L12C2B8-32-TLN
% Mức: TH
\begin{ex}
% ID: L12C2B8-32-TLN
% Mức: TH
Trong dạng “Đặc điểm cấu tạo và chuyển động phân tử ở thể khí”, Một lượng khí có thể tích 4 lít được nén còn 2 lít. Mật độ phân tử tăng bao nhiêu lần?
\shortans{2}
\loigiai{
Số phân tử không đổi nên mật độ tỉ lệ nghịch thể tích: $n_2/n_1=V_1/V_2=2$.
}
\end{ex}

% ===== Câu 126 =====
% ID: L12C2B8-32-TN
% Mức: NB
\dangbt{Giải thích áp suất khí bằng va chạm phân tử}
\begin{ex}
% ID: L12C2B8-32-TN
% Mức: NB
Trong dạng “Giải thích áp suất khí bằng va chạm phân tử”, Kích thước phân tử so với khoảng cách trung bình giữa chúng trong khí lí tưởng được coi là
\choice
{\True rất nhỏ}
{rất lớn}
{bằng nhau}
{không xác định}
\loigiai{
Đáp án A. Trong mô hình khí lí tưởng, kích thước phân tử được bỏ qua so với khoảng cách giữa chúng.
}
\end{ex}

% ===== Câu 127 =====
% ID: L12C2B8-33-TL
% Mức: VD
\dangbt{Đặc điểm cấu tạo và chuyển động phân tử ở thể khí}
\begin{ex}
% ID: L12C2B8-33-TL
% Mức: VD
Trong dạng “Đặc điểm cấu tạo và chuyển động phân tử ở thể khí”, Giải thích vì sao chất khí gây áp suất lên thành bình.
\choiceTF

\loigiai{
Số phân tử khí không đổi nên mật độ phân tử $n = \frac{N}{V}$ tỉ lệ nghịch với thể tích $V$.
Tỉ số mật độ phân tử sau và trước khi nén là: $\frac{n_2}{n_1} = \frac{V_1}{V_2} = \frac{4}{2} = 2$.
Mật độ phân tử tăng 2 lần.

Đáp số: 2
}
\end{ex}

% ===== Câu 128 =====
% ID: L12C2B8-33-TLN
% Mức: NB
\dangbt{Chuyển động Brown và bằng chứng về chuyển động phân tử}
\begin{ex}
% ID: L12C2B8-33-TLN
% Mức: NB
Trong dạng “Chuyển động Brown và bằng chứng về chuyển động phân tử”, kích thước phân tử so với khoảng cách trung bình giữa chúng trong khí lí tưởng được coi là
\choice
{\True rất nhỏ}
{rất lớn}
{bằng nhau}
{không xác định}
\loigiai{
Đáp án A. Trong mô hình khí lí tưởng, kích thước phân tử được bỏ qua so với khoảng cách giữa chúng.
}
\end{ex}

% ===== Câu 129 =====
% ID: L12C2B8-33-TN
% Mức: TH
\dangbt{Giải thích áp suất khí bằng va chạm phân tử}
\begin{ex}
% ID: L12C2B8-33-TN
% Mức: TH
Trong dạng “Giải thích áp suất khí bằng va chạm phân tử”, Lực tương tác giữa các phân tử khí lí tưởng được coi là
\choice
{\True không đáng kể trừ khi va chạm}
{luôn rất lớn}
{lực hút không đổi}
{bằng trọng lực}
\loigiai{
Đáp án A. Ngoài thời gian va chạm rất ngắn, tương tác phân tử được bỏ qua.
}
\end{ex}

% ===== Câu 130 =====
% ID: L12C2B8-34-TL
% Mức: VD
\dangbt{Đặc điểm cấu tạo và chuyển động phân tử ở thể khí}
\begin{ex}
% ID: L12C2B8-34-TL
% Mức: VD
Trong dạng “Đặc điểm cấu tạo và chuyển động phân tử ở thể khí”, So sánh chuyển động phân tử khi nhiệt độ tăng.
\choiceTF

\loigiai{
Nhiệt độ tuyệt đối tăng làm động năng tịnh tiến trung bình và tốc độ đặc trưng của phân tử tăng.
}
\end{ex}

% ===== Câu 131 =====
% ID: L12C2B8-34-TLN
% Mức: VD
\begin{ex}
% ID: L12C2B8-34-TLN
% Mức: VD
Trong dạng “Đặc điểm cấu tạo và chuyển động phân tử ở thể khí”, Khí có $N=1,2\times10^{24}$ phân tử và mật độ $3\times10^{25}\,m^{-3}$. Tính thể tích theo lít.
\shortans{40}
\loigiai{
Động năng tịnh tiến trung bình của phân tử khí lý tưởng tỉ lệ thuận với nhiệt độ tuyệt đối theo công thức $\overline{W_đ} = \frac{3}{2}kT$.
Khi nhiệt độ tăng từ $T_1 = 300 \, \mathrm{K}$ lên $T_2 = 600 \, \mathrm{K}$, tỉ số động năng là $\frac{\overline{W_{đ2}}}{\overline{W_{đ1}}} = \frac{T_2}{T_1} = \frac{600}{300} = 2$.
Vậy động năng tịnh tiến trung bình tăng 2 lần.

Đáp số: 2
}
\end{ex}

% ===== Câu 132 =====
% ID: L12C2B8-34-TN
% Mức: TH
\dangbt{Giải thích áp suất khí bằng va chạm phân tử}
\begin{ex}
% ID: L12C2B8-34-TN
% Mức: TH
Trong dạng “Giải thích áp suất khí bằng va chạm phân tử”, Chuyển động Brown là chuyển động
\choice
{\True hỗn loạn của hạt nhỏ lơ lửng do va chạm phân tử}
{thẳng đều của phân tử khí}
{rơi tự do của hạt bụi}
{tuần hoàn của chất lỏng}
\loigiai{
Đáp án A. Các phân tử môi trường va chạm không cân bằng làm hạt nhỏ chuyển động Brown.
}
\end{ex}

% ===== Câu 133 =====
% ID: L12C2B8-35-TL
% Mức: VD
\dangbt{Đặc điểm cấu tạo và chuyển động phân tử ở thể khí}
\begin{ex}
% ID: L12C2B8-35-TL
% Mức: VD
Trong dạng “Đặc điểm cấu tạo và chuyển động phân tử ở thể khí”, Giải thích hiện tượng mùi nước hoa lan khắp phòng.
\choiceTF

\loigiai{
Phân tử nước hoa bay hơi rồi chuyển động nhiệt hỗn loạn, va chạm và khuếch tán từ nơi nồng độ cao sang thấp.
}
\end{ex}

% ===== Câu 134 =====
% ID: L12C2B8-35-TLN
% Mức: VD
\begin{ex}
% ID: L12C2B8-35-TLN
% Mức: VD
Trong dạng “Đặc điểm cấu tạo và chuyển động phân tử ở thể khí”, Có 2 mol khí. Lấy $N_A=6,02\times10^{23}$. Số phân tử theo đơn vị $10^{23}$ là bao nhiêu?
\shortans{12,04}
\loigiai{
$N=2N_A=12,04\times10^{23}$.
}
\end{ex}

% ===== Câu 135 =====
% ID: L12C2B8-35-TN
% Mức: TH
\dangbt{Giải thích áp suất khí bằng va chạm phân tử}
\begin{ex}
% ID: L12C2B8-35-TN
% Mức: TH
Trong dạng “Giải thích áp suất khí bằng va chạm phân tử”, Áp suất khí lên thành bình sinh ra chủ yếu do
\choice
{\True phân tử va chạm thành bình}
{trọng lượng bình}
{lực hút của Trái Đất}
{thành bình nóng lên}
\loigiai{
Thể tích của khối khí được tính theo công thức $V = \frac{N}{n_0}$.
Thay số: $V = \frac{1,2 \times 10^{24}}{3 \times 10^{25}} = 0,04 \, \mathrm{m^3}$.
Đổi đơn vị: $0,04 \, \mathrm{m^3} = 40 \, \mathrm{lít}$.

Đáp số: 40 lít
}
\end{ex}

% ===== Câu 136 =====
% ID: L12C2B8-36-TL
% Mức: VDC
\dangbt{Đặc điểm cấu tạo và chuyển động phân tử ở thể khí}
\begin{ex}
% ID: L12C2B8-36-TL
% Mức: VDC
Trong dạng “Đặc điểm cấu tạo và chuyển động phân tử ở thể khí”, Nêu giới hạn của mô hình khí lí tưởng.
\choiceTF

\loigiai{
Mô hình phù hợp khi khí loãng, áp suất không quá cao và nhiệt độ không gần vùng hóa lỏng; khi đó thể tích riêng và tương tác phân tử có thể bỏ qua.
}
\end{ex}

% ===== Câu 137 =====
% ID: L12C2B8-36-TLN
% Mức: VDC
\begin{ex}
% ID: L12C2B8-36-TLN
% Mức: VDC
Trong dạng “Đặc điểm cấu tạo và chuyển động phân tử ở thể khí”, Khí chiếm 5 lít, giãn nở thành 15 lít. Mật độ phân tử giảm bao nhiêu lần?
\shortans{3}
\loigiai{
$n_1/n_2=V_2/V_1=15/5=3$.
}
\end{ex}

% ===== Câu 138 =====
% ID: L12C2B8-36-TN
% Mức: VD
\dangbt{Giải thích áp suất khí bằng va chạm phân tử}
\begin{ex}
% ID: L12C2B8-36-TN
% Mức: VD
Trong dạng “Giải thích áp suất khí bằng va chạm phân tử”, Khi nhiệt độ khí tăng, chuyển động nhiệt của phân tử nhìn chung
\choice
{\True nhanh hơn}
{chậm hơn}
{không đổi}
{dừng lại}
\loigiai{
Đáp án A. Nhiệt độ tăng làm động năng tịnh tiến trung bình tăng.
}
\end{ex}

% ===== Câu 139 =====
% ID: L12C2B8-37-TN
% Mức: VD
\begin{ex}
% ID: L12C2B8-37-TN
% Mức: VD
Trong dạng “Giải thích áp suất khí bằng va chạm phân tử”, Ba thông số trạng thái cơ bản của một lượng khí là
\choice
{\True áp suất, thể tích, nhiệt độ}
{khối lượng, trọng lượng, nhiệt lượng}
{vận tốc, gia tốc, thời gian}
{lực, công, công suất}
\loigiai{
Số mol khí $n = 2 \, \mathrm{mol}$, số Avogadro $N_A = 6,02 \times 10^{23} \, \mathrm{phân tử/mol}$.
Số phân tử khí có trong bình là $N = n \cdot N_A = 2 \cdot 6,02 \times 10^{23} = 12,04 \times 10^{23} \, \mathrm{phân tử}$.
Vậy số phân tử theo đơn vị $10^{23}$ là 12,04.

Đáp số: 12,04
}
\end{ex}

% ===== Câu 140 =====
% ID: L12C2B8-38-TN
% Mức: VD
\begin{ex}
% ID: L12C2B8-38-TN
% Mức: VD
Trong dạng “Giải thích áp suất khí bằng va chạm phân tử”, Va chạm giữa các phân tử khí lí tưởng được xem là
\choice
{\True đàn hồi}
{mềm hoàn toàn}
{không bảo toàn động lượng}
{không xảy ra}
\loigiai{
Đáp án A. Mô hình khí lí tưởng coi các va chạm là đàn hồi.
}
\end{ex}

% ===== Câu 141 =====
% ID: L12C2B8-39-TN
% Mức: VDC
\begin{ex}
% ID: L12C2B8-39-TN
% Mức: VDC
Trong dạng “Giải thích áp suất khí bằng va chạm phân tử”, Khi nén khí ở nhiệt độ không đổi, mật độ phân tử
\choice
{\True tăng}
{giảm}
{không đổi}
{bằng không}
\loigiai{
Đáp án A. Cùng số phân tử trong thể tích nhỏ hơn nên mật độ phân tử tăng.
}
\end{ex}

% ===== Câu 142 =====
% ID: L12C2B8-40-TN
% Mức: VDC
\dangbt{Chuyển động Brown và bằng chứng về chuyển động phân tử}
\begin{ex}
% ID: L12C2B8-40-TN
% Mức: VDC
Xét các phát biểu về mô hình khí lí tưởng và giới hạn áp dụng của nó:
\choiceTF
{\True Mô hình khí lí tưởng coi các phân tử khí là các chất điểm và chỉ tương tác với nhau khi va chạm.}
{Mô hình khí lí tưởng áp dụng chính xác nhất cho mọi chất khí ở mọi điều kiện nhiệt độ và áp suất.}
{\True Khi áp suất chất khí rất cao, thể tích riêng của các phân tử không còn có thể bỏ qua so với thể tích bình chứa.}
{\True Ở nhiệt độ rất thấp, gần nhiệt độ hóa lỏng, lực tương tác giữa các phân tử khí trở nên đáng kể, làm mô hình khí lí tưởng không còn chính xác.}
\loigiai{
\begin{itemchoice}
\item (Đúng) Mô hình khí lí tưởng coi các phân tử khí là các chất điểm và chỉ tương tác với nhau khi va chạm. (Vì): Theo giả thuyết của mô hình khí lí tưởng, các phân tử khí được coi là chất điểm và chỉ tương tác với nhau khi va chạm đàn hồi
\item (Sai) Mô hình khí lí tưởng áp dụng chính xác nhất cho mọi chất khí ở mọi điều kiện nhiệt độ và áp suất. (Vì): Mô hình khí lí tưởng chỉ là mô hình gần đúng, nó không áp dụng chính xác cho mọi chất khí ở mọi điều kiện, đặc biệt là khi áp suất rất cao hoặc nhiệt độ rất thấp
\item (Đúng) Khi áp suất chất khí rất cao, thể tích riêng của các phân tử không còn có thể bỏ qua so với thể tích bình chứa. (Vì): Khi áp suất chất khí rất cao, khoảng cách giữa các phân tử giảm mạnh, lúc này thể tích riêng của các phân tử không còn đủ nhỏ để có thể bỏ qua so với thể tích của bình chứa
\item (Đúng) Ở nhiệt độ rất thấp, gần nhiệt độ hóa lỏng, lực tương tác giữa các phân tử khí trở nên đáng kể, làm mô hình khí lí tưởng không còn chính xác. (Vì): Ở nhiệt độ rất thấp, động năng của các phân tử giảm, lực tương tác giữa các phân tử trở nên đáng kể, khiến chất khí không còn tuân theo các định luật của khí lí tưởng.

Đáp án: A=Đúng; B=Sai; C=Đúng; D=Đúng
\end{itemchoice}
}
\end{ex}

% ===== Câu 143 =====
% ID: L12C2B8-41-TN
% Mức: NB
\dangbt{Vận dụng mô hình động học phân tử vào thực tế}
\begin{ex}
% ID: L12C2B8-41-TN
% Mức: NB
Trong dạng “Vận dụng mô hình động học phân tử vào thực tế”, Theo mô hình động học phân tử, các phân tử khí
\choice
{\True chuyển động hỗn loạn không ngừng}
{đứng yên tại vị trí cân bằng}
{chỉ chuyển động khi được nung nóng}
{chuyển động cùng một hướng}
\loigiai{
Thể tích ban đầu $V_1 = 5 \, \mathrm{lít}$, thể tích sau khi giãn nở $V_2 = 15 \, \mathrm{lít}$.
Mật độ phân tử tỉ lệ nghịch với thể tích: $n = \frac{N}{V} \Rightarrow \frac{n_1}{n_2} = \frac{V_2}{V_1}$.
Thay số: $\frac{n_1}{n_2} = \frac{15}{5} = 3$.
Mật độ phân tử giảm 3 lần.
Đáp số: 3
}
\end{ex}

% ===== Câu 144 =====
% ID: L12C2B8-42-TN
% Mức: NB
\begin{ex}
% ID: L12C2B8-42-TN
% Mức: NB
Trong dạng “Vận dụng mô hình động học phân tử vào thực tế”, Kích thước phân tử so với khoảng cách trung bình giữa chúng trong khí lí tưởng được coi là
\choice
{\True rất nhỏ}
{rất lớn}
{bằng nhau}
{không xác định}
\loigiai{
Đáp án A. Trong mô hình khí lí tưởng, kích thước phân tử được bỏ qua so với khoảng cách giữa chúng.
}
\end{ex}

% ===== Câu 145 =====
% ID: L12C2B8-43-TN
% Mức: TH
\begin{ex}
% ID: L12C2B8-43-TN
% Mức: TH
Trong dạng “Vận dụng mô hình động học phân tử vào thực tế”, Lực tương tác giữa các phân tử khí lí tưởng được coi là
\choice
{\True không đáng kể trừ khi va chạm}
{luôn rất lớn}
{lực hút không đổi}
{bằng trọng lực}
\loigiai{
Đáp án A. Ngoài thời gian va chạm rất ngắn, tương tác phân tử được bỏ qua.
}
\end{ex}

% ===== Câu 146 =====
% ID: L12C2B8-44-TN
% Mức: TH
\begin{ex}
% ID: L12C2B8-44-TN
% Mức: TH
Trong dạng “Vận dụng mô hình động học phân tử vào thực tế”, Chuyển động Brown là chuyển động
\choice
{\True hỗn loạn của hạt nhỏ lơ lửng do va chạm phân tử}
{thẳng đều của phân tử khí}
{rơi tự do của hạt bụi}
{tuần hoàn của chất lỏng}
\loigiai{
Đáp án A. Các phân tử môi trường va chạm không cân bằng làm hạt nhỏ chuyển động Brown.
}
\end{ex}

% ===== Câu 147 =====
% ID: L12C2B8-45-TN
% Mức: TH
\begin{ex}
% ID: L12C2B8-45-TN
% Mức: TH
Trong dạng “Vận dụng mô hình động học phân tử vào thực tế”, Áp suất khí lên thành bình sinh ra chủ yếu do
\choice
{\True phân tử va chạm thành bình}
{trọng lượng bình}
{lực hút của Trái Đất}
{thành bình nóng lên}
\loigiai{
Đáp án A. Xung lượng truyền trong các va chạm phân tử với thành bình tạo áp suất.
}
\end{ex}

% ===== Câu 148 =====
% ID: L12C2B8-46-TN
% Mức: VD
\begin{ex}
% ID: L12C2B8-46-TN
% Mức: VD
Trong dạng “Vận dụng mô hình động học phân tử vào thực tế”, Khi nhiệt độ khí tăng, chuyển động nhiệt của phân tử nhìn chung
\choice
{\True nhanh hơn}
{chậm hơn}
{không đổi}
{dừng lại}
\loigiai{
Đáp án A. Nhiệt độ tăng làm động năng tịnh tiến trung bình tăng.
}
\end{ex}

% ===== Câu 149 =====
% ID: L12C2B8-47-TN
% Mức: VD
\begin{ex}
% ID: L12C2B8-47-TN
% Mức: VD
Trong dạng “Vận dụng mô hình động học phân tử vào thực tế”, Ba thông số trạng thái cơ bản của một lượng khí là
\choice
{\True áp suất, thể tích, nhiệt độ}
{khối lượng, trọng lượng, nhiệt lượng}
{vận tốc, gia tốc, thời gian}
{lực, công, công suất}
\loigiai{
Đáp án A. Trạng thái lượng khí xác định bởi p, $V$, T.
}
\end{ex}

% ===== Câu 150 =====
% ID: L12C2B8-48-TN
% Mức: VD
\begin{ex}
% ID: L12C2B8-48-TN
% Mức: VD
Trong dạng “Vận dụng mô hình động học phân tử vào thực tế”, Va chạm giữa các phân tử khí lí tưởng được xem là
\choice
{\True đàn hồi}
{mềm hoàn toàn}
{không bảo toàn động lượng}
{không xảy ra}
\loigiai{
Đáp án A. Mô hình khí lí tưởng coi các va chạm là đàn hồi.
}
\end{ex}

% ===== Câu 151 =====
% ID: L12C2B8-49-TN
% Mức: VDC
\begin{ex}
% ID: L12C2B8-49-TN
% Mức: VDC
Trong dạng “Vận dụng mô hình động học phân tử vào thực tế”, Khi nén khí ở nhiệt độ không đổi, mật độ phân tử
\choice
{\True tăng}
{giảm}
{không đổi}
{bằng không}
\loigiai{
Đáp án A. Cùng số phân tử trong thể tích nhỏ hơn nên mật độ phân tử tăng.
}
\end{ex}

% ===== Câu 152 =====
% ID: L12C2B8-50-TN
% Mức: VDC
\begin{ex}
% ID: L12C2B8-50-TN
% Mức: VDC
Trong dạng “Vận dụng mô hình động học phân tử vào thực tế”, Hiện tượng khuếch tán chứng tỏ các phân tử
\choice
{\True chuyển động không ngừng}
{có kích thước vĩ mô}
{đứng yên}
{không tương tác khi va chạm}
\loigiai{
Đáp án A. Khuếch tán là hệ quả trực tiếp của chuyển động nhiệt hỗn loạn.
}
\end{ex}

% ===== Câu 153 =====
% ID: L12C2B8-51-TN
% Mức: NB
\dangbt{Đặc điểm cấu tạo và chuyển động phân tử ở thể khí}
\begin{ex}
% ID: L12C2B8-51-TN
% Mức: NB
Trong dạng “Đặc điểm cấu tạo và chuyển động phân tử ở thể khí”, Theo mô hình động học phân tử, các phân tử khí
\choice
{\True chuyển động hỗn loạn không ngừng}
{đứng yên tại vị trí cân bằng}
{chỉ chuyển động khi được nung nóng}
{chuyển động cùng một hướng}
\loigiai{
Đáp án A. Phân tử khí luôn chuyển động nhiệt hỗn loạn không ngừng.
}
\end{ex}

% ===== Câu 154 =====
% ID: L12C2B8-52-TN
% Mức: NB
\begin{ex}
% ID: L12C2B8-52-TN
% Mức: NB
Trong dạng “Đặc điểm cấu tạo và chuyển động phân tử ở thể khí”, Kích thước phân tử so với khoảng cách trung bình giữa chúng trong khí lí tưởng được coi là
\choice
{\True rất nhỏ}
{rất lớn}
{bằng nhau}
{không xác định}
\loigiai{
Đáp án A. Trong mô hình khí lí tưởng, kích thước phân tử được bỏ qua so với khoảng cách giữa chúng.
}
\end{ex}

% ===== Câu 155 =====
% ID: L12C2B8-53-TN
% Mức: TH
\begin{ex}
% ID: L12C2B8-53-TN
% Mức: TH
Trong dạng “Đặc điểm cấu tạo và chuyển động phân tử ở thể khí”, Lực tương tác giữa các phân tử khí lí tưởng được coi là
\choice
{\True không đáng kể trừ khi va chạm}
{luôn rất lớn}
{lực hút không đổi}
{bằng trọng lực}
\loigiai{
Đáp án A. Ngoài thời gian va chạm rất ngắn, tương tác phân tử được bỏ qua.
}
\end{ex}

% ===== Câu 156 =====
% ID: L12C2B8-54-TN
% Mức: TH
\begin{ex}
% ID: L12C2B8-54-TN
% Mức: TH
Trong dạng “Đặc điểm cấu tạo và chuyển động phân tử ở thể khí”, Chuyển động Brown là chuyển động
\choice
{\True hỗn loạn của hạt nhỏ lơ lửng do va chạm phân tử}
{thẳng đều của phân tử khí}
{rơi tự do của hạt bụi}
{tuần hoàn của chất lỏng}
\loigiai{
Đáp án A. Các phân tử môi trường va chạm không cân bằng làm hạt nhỏ chuyển động Brown.
}
\end{ex}

% ===== Câu 157 =====
% ID: L12C2B8-55-TN
% Mức: TH
\begin{ex}
% ID: L12C2B8-55-TN
% Mức: TH
Trong dạng “Đặc điểm cấu tạo và chuyển động phân tử ở thể khí”, Áp suất khí lên thành bình sinh ra chủ yếu do
\choice
{\True phân tử va chạm thành bình}
{trọng lượng bình}
{lực hút của Trái Đất}
{thành bình nóng lên}
\loigiai{
Đáp án A. Xung lượng truyền trong các va chạm phân tử với thành bình tạo áp suất.
}
\end{ex}

% ===== Câu 158 =====
% ID: L12C2B8-56-TN
% Mức: VD
\begin{ex}
% ID: L12C2B8-56-TN
% Mức: VD
Trong dạng “Đặc điểm cấu tạo và chuyển động phân tử ở thể khí”, Khi nhiệt độ khí tăng, chuyển động nhiệt của phân tử nhìn chung
\choice
{\True nhanh hơn}
{chậm hơn}
{không đổi}
{dừng lại}
\loigiai{
Đáp án A. Nhiệt độ tăng làm động năng tịnh tiến trung bình tăng.
}
\end{ex}

% ===== Câu 159 =====
% ID: L12C2B8-57-TN
% Mức: VD
\begin{ex}
% ID: L12C2B8-57-TN
% Mức: VD
Trong dạng “Đặc điểm cấu tạo và chuyển động phân tử ở thể khí”, Ba thông số trạng thái cơ bản của một lượng khí là
\choice
{\True áp suất, thể tích, nhiệt độ}
{khối lượng, trọng lượng, nhiệt lượng}
{vận tốc, gia tốc, thời gian}
{lực, công, công suất}
\loigiai{
Đáp án A. Trạng thái lượng khí xác định bởi p, $V$, T.
}
\end{ex}

% ===== Câu 160 =====
% ID: L12C2B8-58-TN
% Mức: VD
\begin{ex}
% ID: L12C2B8-58-TN
% Mức: VD
Trong dạng “Đặc điểm cấu tạo và chuyển động phân tử ở thể khí”, Va chạm giữa các phân tử khí lí tưởng được xem là
\choice
{\True đàn hồi}
{mềm hoàn toàn}
{không bảo toàn động lượng}
{không xảy ra}
\loigiai{
Đáp án A. Mô hình khí lí tưởng coi các va chạm là đàn hồi.
}
\end{ex}

% ===== Câu 161 =====
% ID: L12C2B8-59-TN
% Mức: VDC
\begin{ex}
% ID: L12C2B8-59-TN
% Mức: VDC
Trong dạng “Đặc điểm cấu tạo và chuyển động phân tử ở thể khí”, Khi nén khí ở nhiệt độ không đổi, mật độ phân tử
\choice
{\True tăng}
{giảm}
{không đổi}
{bằng không}
\loigiai{
Đáp án A. Cùng số phân tử trong thể tích nhỏ hơn nên mật độ phân tử tăng.
}
\end{ex}

% ===== Câu 162 =====
% ID: L12C2B8-60-TN
% Mức: VDC
\begin{ex}
% ID: L12C2B8-60-TN
% Mức: VDC
Trong dạng “Đặc điểm cấu tạo và chuyển động phân tử ở thể khí”, Hiện tượng khuếch tán chứng tỏ các phân tử
\choice
{\True chuyển động không ngừng}
{có kích thước vĩ mô}
{đứng yên}
{không tương tác khi va chạm}
\loigiai{
Đáp án A. Khuếch tán là hệ quả trực tiếp của chuyển động nhiệt hỗn loạn.
}
\end{ex}
